%% 6-14 微分方程式がいくつも：連立微分方程式（第14回：組 V6-8。ページ制限なし。偶数ページ）
%% ほかの巻・物理の本との約束：
%%   行列（第7巻）は使わない。係数を「表の形」に並べて見せるだけにし、7-8・7-9（固有ベクトル・対角化）、9-10（状態方程式）を予告する。
%%   2階の式は「位置と速度」の1階の式2つにできる（9-10 の状態の考えの入口）。表計算の型は 6-2（速度を先に進める）。
%%   和と差に分ける（同じ水そう2つ：合計は一定、差は 6-4 の形）、振動モード（いっしょ・反対）、消去法（3-6 の代入法と同じ）。
%%   直列の水そうは 6-12 の重解 te^{st} と同じ形、親と子の放射性物質（ベイトマン 1910年、原-12）。熱・電気も同じ形（6-15 の相似へ）。
%%   コラム：カオス（ローレンツ 1961年・1963年の論文）。
\makeatletter
% 小見出し（太字の1行）  \kviohead{見出し}
\newcommand\kviohead[1]{\par\addvspace{6pt}\Needspace{6\baselineskip}\noindent{\sffamily\bfseries\color{\bkTT}#1}\par\nobreak\vspace{2pt}}
% 囲みの中の小見出し（それだけで1行）
\newcommand\kviobox[1]{\par\addvspace{3pt}\noindent\textbf{#1}\par\vspace{3pt}\noindent\ignorespaces}
% 式を、まん中に1行で置く
\newcommand\kvioline[1]{\par\addvspace{3pt}{\centering #1\par}\addvspace{3pt}\noindent\ignorespaces}
% 大事な式のわく  \kvioframe[はば]{中身}
\newcommand\kvioframe[2][100mm]{\par\addvspace{4pt}{\centering\tikz\node[draw=\bkTT, line width=0.9pt, fill=bkTheme!10, rounded corners=1.5mm, inner xsep=6mm, inner ysep=2.5mm, align=center]{#2};\par}\addvspace{4pt}\noindent\ignorespaces}
% 例題の途中式の2行め以降
\newcommand\kvioa{\par\vspace{2pt}\noindent\rightskip=0pt plus 8em\relax\hangindent=2.6\zw\hangafter=0\relax}
\newcommand\kviob{\par\vspace{2pt}\noindent\rightskip=0pt plus 8em\relax\hangindent=4.3\zw\hangafter=0\relax}
\newcommand\kvioq{\par\vspace{2pt}\noindent\rightskip=0pt plus 8em\relax\hangindent=4.3\zw\hangafter=1\hspace*{2.6\zw}}
% 例題の式の段の1行め（行の右をそろえない）
\newcommand\kvios{\rightskip=0pt plus 8em\relax}
% 水そう  \kviotank{(x,y)}{はば}{高さ}{水の深さ}
\newcommand\kviotank[4]{\begin{scope}[shift={#1}]
  \fill[bkBlue!22] (0,0) rectangle (#2,#4);
  \draw[bkBlue!70!black, line width=0.5pt] (0,#4) -- (#2,#4);
  \draw[bkInk!75, line width=0.9pt] (0,#3) -- (0,0) -- (#2,0) -- (#2,#3);
\end{scope}}
% ばね（横）  \kviospring{(左のはし)}{(右のはし)}
\newcommand\kviospring[2]{\draw[bkInk, line width=0.8pt, decorate, decoration={zigzag, amplitude=1.3mm, segment length=2.2mm, pre length=1.5mm, post length=1.5mm}] #1 -- #2;}
% 台車（おもり）  \kviocart{(左下)}{文字}{色}
\newcommand\kviocart[3]{\begin{scope}[shift={#1}]
  \filldraw[fill=#3!25, draw=#3!80!black, line width=0.8pt, rounded corners=0.8mm] (0,1.6) rectangle (12,10);
  \filldraw[fill=bkInk!60, draw=bkInk] (3,0.8) circle (0.8) (9,0.8) circle (0.8);
  \node[font=\sffamily\bfseries\small] at (6,5.8) {#2};
\end{scope}}
% かべ  \kviowall{(x,y)}{高さ}{向き：1 なら右に面、-1 なら左に面}
\newcommand\kviowall[3]{\begin{scope}[shift={#1}]
  \draw[bkInk, line width=0.9pt] (0,0) -- (0,#2);
  \foreach \kvioy in {0,2.5,...,#2} {\draw[bkInk!60, line width=0.4pt] (0,\kvioy) -- ({-1.6*#3},{\kvioy+1.6});}
\end{scope}}
\makeatother
%% グラフの点（Python の RK4 で計算した。組 V6-8）
\def\kviolorenzxz{(-9.2,29.3) (-8.8,29.0) (-8.3,28.6) (-8.0,28.0) (-7.6,27.3) (-7.3,26.6) (-7.2,26.0) (-7.1,25.4) (-7.2,24.9) (-7.3,24.6) (-7.5,24.5) (-7.8,24.5) (-8.2,24.8) (-8.6,25.3) (-9.0,25.9) (-9.4,26.7) (-9.7,27.5) (-9.8,28.3) (-9.9,28.9) (-9.7,29.4) (-9.4,29.5) (-9.1,29.4) (-8.6,29.0) (-8.2,28.5) (-7.8,27.8) (-7.4,27.1) (-7.2,26.4) (-7.1,25.7) (-7.0,25.1) (-7.1,24.7) (-7.3,24.4) (-7.6,24.3) (-7.9,24.4) (-8.3,24.8) (-8.7,25.3) (-9.2,26.0) (-9.5,26.9) (-9.8,27.8) (-9.9,28.6) (-9.9,29.2) (-9.7,29.6) (-9.4,29.7) (-9.0,29.4) (-8.5,29.0) (-8.0,28.4) (-7.6,27.6) (-7.3,26.9) (-7.1,26.1) (-6.9,25.4) (-6.9,24.9) (-7.0,24.4) (-7.3,24.2) (-7.6,24.1) (-8.0,24.3) (-8.4,24.7) (-8.9,25.4) (-9.3,26.2) (-9.7,27.1) (-9.9,28.1) (-10.0,28.9) (-10.0,29.5) (-9.7,29.8) (-9.3,29.8) (-8.9,29.5) (-8.4,28.9) (-7.9,28.2) (-7.5,27.4) (-7.1,26.6) (-6.9,25.9) (-6.8,25.2) (-6.8,24.6) (-7.0,24.2) (-7.3,24.0) (-7.6,24.0) (-8.0,24.2) (-8.5,24.7) (-9.0,25.4) (-9.5,26.4) (-9.8,27.4) (-10.1,28.4) (-10.1,29.2) (-10.0,29.8) (-9.7,30.0) (-9.3,29.9) (-8.8,29.5) (-8.2,28.9) (-7.7,28.1) (-7.3,27.2) (-7.0,26.4) (-6.8,25.6) (-6.7,24.9) (-6.8,24.3) (-6.9,23.9) (-7.2,23.7) (-7.6,23.8) (-8.1,24.1) (-8.6,24.7) (-9.1,25.5) (-9.6,26.5) (-10.0,27.6) (-10.2,28.7) (-10.2,29.5) (-10.1,30.1) (-9.7,30.2) (-9.2,30.0) (-8.6,29.5) (-8.1,28.8) (-7.5,28.0) (-7.1,27.1) (-6.8,26.2) (-6.6,25.3) (-6.6,24.6) (-6.7,24.0) (-6.9,23.6) (-7.2,23.5) (-7.7,23.6) (-8.2,24.0) (-8.7,24.6) (-9.3,25.5) (-9.8,26.7) (-10.1,27.8) (-10.4,28.9) (-10.3,29.8) (-10.1,30.3) (-9.7,30.4) (-9.2,30.2) (-8.5,29.6) (-7.9,28.8) (-7.4,27.8) (-6.9,26.9) (-6.6,25.9) (-6.5,25.1) (-6.5,24.3) (-6.6,23.8) (-6.8,23.4) (-7.2,23.2) (-7.7,23.4) (-8.2,23.8) (-8.8,24.5) (-9.4,25.6) (-9.9,26.8) (-10.3,28.0) (-10.5,29.2) (-10.5,30.1) (-10.2,30.6) (-9.7,30.7) (-9.1,30.3) (-8.4,29.6) (-7.8,28.7) (-7.2,27.7) (-6.8,26.7) (-6.5,25.7) (-6.3,24.8) (-6.3,24.1) (-6.4,23.5) (-6.7,23.1) (-7.1,23.0) (-7.6,23.1) (-8.2,23.6) (-8.9,24.4) (-9.5,25.5) (-10.0,26.8) (-10.4,28.2) (-10.6,29.4) (-10.6,30.4) (-10.3,30.9) (-9.7,30.9) (-9.1,30.5) (-8.4,29.7) (-7.7,28.8) (-7.1,27.7) (-6.6,26.6) (-6.3,25.6) (-6.1,24.6) (-6.2,23.8) (-6.3,23.2) (-6.6,22.8) (-7.0,22.6) (-7.6,22.8) (-8.2,23.3) (-8.9,24.2) (-9.5,25.4) (-10.1,26.8) (-10.6,28.3) (-10.8,29.6) (-10.7,30.6) (-10.4,31.2) (-9.8,31.2) (-9.1,30.7) (-8.3,29.9) (-7.6,28.8) (-6.9,27.7) (-6.5,26.5) (-6.1,25.4) (-6.0,24.4) (-6.0,23.6) (-6.1,22.9) (-6.4,22.5) (-6.9,22.3) (-7.5,22.5) (-8.1,23.0) (-8.8,23.9) (-9.6,25.1) (-10.2,26.7) (-10.7,28.3) (-11.0,29.8) (-10.9,30.9) (-10.5,31.5) (-9.9,31.5) (-9.2,31.0) (-8.3,30.1) (-7.5,29.0) (-6.8,27.8) (-6.3,26.5) (-6.0,25.4) (-5.8,24.3) (-5.8,23.4) (-5.9,22.6) (-6.2,22.1) (-6.7,21.9) (-7.3,22.0) (-8.0,22.5) (-8.7,23.4) (-9.5,24.7) (-10.3,26.3) (-10.8,28.1) (-11.1,29.8) (-11.1,31.1) (-10.8,31.8) (-10.1,31.8) (-9.3,31.3) (-8.4,30.4) (-7.5,29.2) (-6.8,27.9) (-6.2,26.6) (-5.8,25.4) (-5.6,24.2) (-5.5,23.2) (-5.7,22.4) (-6.0,21.8) (-6.4,21.5) (-7.0,21.5) (-7.7,21.9) (-8.5,22.8) (-9.4,24.1) (-10.2,25.8) (-10.9,27.7) (-11.3,29.6) (-11.4,31.1) (-11.1,32.0) (-10.4,32.3) (-9.6,31.8) (-8.6,30.9) (-7.6,29.7) (-6.8,28.3) (-6.1,26.9) (-5.6,25.6) (-5.4,24.3) (-5.3,23.2) (-5.3,22.3) (-5.6,21.6) (-6.0,21.1) (-6.6,21.0) (-7.3,21.2) (-8.2,22.0) (-9.1,23.2) (-10.0,24.8) (-10.8,26.9) (-11.4,29.0) (-11.6,30.9) (-11.5,32.2) (-10.9,32.7) (-10.0,32.5) (-9.0,31.6) (-7.9,30.4) (-6.9,28.9) (-6.1,27.4) (-5.5,26.0) (-5.2,24.6) (-5.0,23.3) (-5.0,22.3) (-5.2,21.4) (-5.5,20.8) (-6.1,20.4) (-6.8,20.5) (-7.6,20.9) (-8.5,21.9) (-9.5,23.5) (-10.5,25.5) (-11.3,27.8) (-11.8,30.1) (-11.9,32.0) (-11.5,33.0) (-10.7,33.2) (-9.6,32.6) (-8.4,31.4) (-7.3,29.9) (-6.3,28.3) (-5.6,26.7) (-5.0,25.2) (-4.7,23.8) (-4.6,22.5) (-4.7,21.4) (-5.0,20.6) (-5.4,20.0) (-6.0,19.7) (-6.8,19.8) (-7.7,20.5) (-8.7,21.6) (-9.8,23.4) (-10.9,25.8) (-11.7,28.4) (-12.2,30.9) (-12.2,32.8) (-11.6,33.8) (-10.6,33.7) (-9.4,32.8) (-8.1,31.4) (-6.9,29.7) (-5.9,27.9) (-5.1,26.2) (-4.6,24.7) (-4.3,23.2) (-4.2,21.9) (-4.3,20.8) (-4.6,19.8) (-5.0,19.2) (-5.7,18.9) (-6.5,18.9) (-7.4,19.5) (-8.6,20.8) (-9.8,22.7) (-11.0,25.2) (-12.0,28.2) (-12.6,31.1) (-12.6,33.3) (-12.0,34.5) (-10.9,34.4) (-9.5,33.4) (-8.1,31.8) (-6.7,30.0) (-5.6,28.1) (-4.7,26.3) (-4.1,24.6) (-3.8,23.0) (-3.7,21.6) (-3.7,20.3) (-3.9,19.3) (-4.4,18.5) (-4.9,17.9) (-5.7,17.7) (-6.6,18.0) (-7.8,18.9) (-9.1,20.6) (-10.4,23.0) (-11.7,26.2) (-12.7,29.7) (-13.1,32.8) (-12.9,34.9) (-11.9,35.5) (-10.5,34.9) (-8.8,33.3) (-7.2,31.4) (-5.8,29.3) (-4.7,27.3) (-3.9,25.4) (-3.3,23.6) (-3.0,22.0) (-3.0,20.5) (-3.0,19.2) (-3.3,18.1) (-3.7,17.2) (-4.2,16.5) (-5.0,16.2) (-5.9,16.3) (-7.1,17.0) (-8.5,18.5) (-10.0,20.9) (-11.6,24.3) (-12.9,28.4) (-13.7,32.4) (-13.7,35.4) (-12.9,36.7) (-11.3,36.3) (-9.4,34.7) (-7.5,32.6) (-5.8,30.2) (-4.4,28.0) (-3.4,25.9) (-2.7,23.9) (-2.2,22.2) (-2.0,20.5) (-1.9,19.1) (-2.0,17.7) (-2.2,16.5) (-2.5,15.5) (-2.9,14.6) (-3.5,14.0) (-4.3,13.6) (-5.3,13.7) (-6.6,14.4) (-8.2,16.0) (-10.0,18.7) (-11.9,22.7) (-13.6,27.8) (-14.7,33.0) (-14.8,37.1) (-13.8,38.8) (-11.9,38.2) (-9.5,36.1) (-7.1,33.4) (-5.0,30.7) (-3.4,28.2) (-2.2,25.9) (-1.3,23.8) (-0.8,22.0) (-0.4,20.3) (-0.1,18.7) (0.1,17.3) (0.2,16.0) (0.3,14.7) (0.4,13.6) (0.5,12.6) (0.7,11.6) (0.9,10.8) (1.1,10.0) (1.4,9.3) (1.9,8.7) (2.4,8.3) (3.2,8.1) (4.2,8.3) (5.6,9.0) (7.3,10.5) (9.5,13.5) (12.0,18.5) (14.6,25.4) (16.5,33.5) (17.1,40.1) (15.8,42.9) (13.0,41.9) (9.6,38.6) (6.3,34.9) (3.5,31.6) (1.3,28.9) (-0.2,26.6) (-1.3,24.6) (-2.2,23.0) (-2.8,21.6) (-3.4,20.4) (-4.1,19.4) (-4.7,18.8) (-5.5,18.5) (-6.4,18.6) (-7.5,19.3) (-8.7,20.7) (-10.0,22.8) (-11.3,25.6) (-12.2,28.8) (-12.8,31.8) (-12.7,33.9) (-12.0,34.9) (-10.8,34.6) (-9.3,33.4) (-7.7,31.7) (-6.4,29.7) (-5.2,27.8) (-4.4,25.9) (-3.8,24.2) (-3.5,22.6) (-3.4,21.2) (-3.5,19.9) (-3.7,18.8) (-4.1,18.0) (-4.7,17.4) (-5.5,17.2) (-6.4,17.4) (-7.6,18.3) (-8.9,19.9) (-10.4,22.4) (-11.8,25.7) (-12.8,29.4) (-13.4,32.8) (-13.2,35.2) (-12.2,36.0) (-10.7,35.4) (-8.9,33.8) (-7.2,31.7) (-5.7,29.5) (-4.5,27.4) (-3.7,25.5) (-3.1,23.6) (-2.7,22.0) (-2.6,20.4) (-2.6,19.1) (-2.8,17.9) (-3.1,16.8) (-3.6,16.0) (-4.3,15.5) (-5.1,15.3) (-6.2,15.6) (-7.5,16.6) (-9.0,18.6) (-10.7,21.5) (-12.3,25.5) (-13.6,30.1) (-14.2,34.2) (-13.9,36.8) (-12.6,37.5) (-10.7,36.4) (-8.6,34.3) (-6.6,31.8) (-4.9,29.4) (-3.6,27.1) (-2.6,25.0) (-2.0,23.1) (-1.6,21.4) (-1.4,19.8) (-1.3,18.3) (-1.3,16.9) (-1.4,15.7) (-1.6,14.6) (-2.0,13.6) (-2.4,12.8) (-3.0,12.2) (-3.7,11.8) (-4.7,11.8) (-5.9,12.3) (-7.5,13.7) (-9.4,16.3) (-11.5,20.3) (-13.6,25.8) (-15.2,32.1) (-15.7,37.3) (-14.9,40.1) (-12.8,39.8) (-10.1,37.6) (-7.3,34.6) (-4.9,31.6) (-3.0,28.9) (-1.6,26.5) (-0.6,24.4) (0.2,22.5) (0.7,20.8) (1.1,19.3) (1.5,17.9) (1.8,16.6) (2.2,15.6) (2.7,14.7) (3.3,14.0) (4.1,13.6) (5.1,13.6) (6.3,14.2) (7.9,15.6) (9.6,18.1) (11.5,21.9) (13.3,26.8) (14.6,32.2) (15.0,36.6) (14.1,38.7) (12.3,38.5) (9.9,36.6) (7.5,33.9) (5.3,31.2) (3.6,28.6) (2.3,26.3) (1.4,24.2) (0.8,22.3) (0.4,20.6) (0.1,19.0) (-0.1,17.5) (-0.3,16.2) (-0.4,15.0) (-0.5,13.8) (-0.7,12.8) (-0.9,11.8) (-1.1,11.0) (-1.4,10.2) (-1.8,9.6) (-2.3,9.1) (-3.1,8.8) (-4.0,8.8) (-5.2,9.3) (-6.9,10.5) (-8.9,13.1) (-11.3,17.3) (-13.8,23.6) (-15.9,31.2) (-16.9,38.2) (-16.2,42.1) (-13.8,42.1) (-10.6,39.4) (-7.3,35.8) (-4.4,32.4) (-2.1,29.5) (-0.4,27.1) (0.8,25.1) (1.7,23.3) (2.3,21.7) (2.9,20.4) (3.5,19.3) (4.1,18.5) (4.8,18.0) (5.7,17.8) (6.7,18.2) (7.9,19.1) (9.2,20.9) (10.6,23.4) (11.9,26.6) (12.8,30.1) (13.1,33.1) (12.8,35.0) (11.8,35.5) (10.3,34.7) (8.6,33.1) (7.0,31.1) (5.7,29.0) (4.6,27.0) (3.8,25.1) (3.3,23.4) (3.0,21.8) (3.0,20.4) (3.1,19.1) (3.3,18.0) (3.7,17.1) (4.3,16.5) (5.1,16.2) (6.1,16.4) (7.3,17.2) (8.7,18.8) (10.2,21.3) (11.8,24.8) (13.1,28.9) (13.8,32.9) (13.7,35.7) (12.7,36.7) (11.1,36.2) (9.1,34.5) (7.2,32.2) (5.6,29.9) (4.3,27.7) (3.3,25.6) (2.6,23.7) (2.2,21.9) (2.0,20.3) (1.9,18.9) (2.0,17.6) (2.2,16.4) (2.6,15.4) (3.0,14.5) (3.7,13.9) (4.5,13.7) (5.6,13.9) (6.9,14.7) (8.5,16.4) (10.3,19.4) (12.2,23.6) (13.8,28.7) (14.8,33.8) (14.7,37.5) (13.5,38.8) (11.5,37.9) (9.0,35.6) (6.7,32.9) (4.7,30.2) (3.2,27.8) (2.0,25.5) (1.3,23.5) (0.7,21.7) (0.4,20.0) (0.1,18.5) (-0.0,17.0) (-0.2,15.7) (-0.3,14.5) (-0.4,13.4) (-0.5,12.4) (-0.6,11.5) (-0.8,10.6) (-1.0,9.8) (-1.3,9.1) (-1.7,8.6) (-2.2,8.1) (-2.9,7.8) (-3.8,7.9) (-5.1,8.4) (-6.7,9.6) (-8.7,12.1) (-11.2,16.4) (-13.8,22.9) (-16.1,30.9) (-17.2,38.4) (-16.5,42.7) (-14.1,42.7) (-10.7,39.9) (-7.2,36.1) (-4.2,32.6) (-1.8,29.7) (-0.1,27.3) (1.2,25.3) (2.1,23.6) (2.8,22.1) (3.4,20.9) (4.0,19.9) (4.7,19.2) (5.5,18.8) (6.4,18.9) (7.4,19.5) (8.6,20.7) (9.8,22.7) (11.0,25.3) (12.0,28.3) (12.6,31.2) (12.6,33.4) (12.0,34.6) (10.9,34.5) (9.5,33.4) (8.0,31.8) (6.7,30.0) (5.5,28.1) (4.7,26.2) (4.1,24.5) (3.7,22.9) (3.6,21.5) (3.7,20.3) (3.9,19.2) (4.3,18.4) (4.9,17.8) (5.6,17.6) (6.6,17.9) (7.7,18.8) (9.0,20.4) (10.4,22.9) (11.7,26.0) (12.7,29.6) (13.2,32.8) (13.0,34.9) (12.0,35.6) (10.6,35.0) (8.9,33.4) (7.2,31.5) (5.8,29.4) (4.7,27.3) (3.8,25.4) (3.3,23.6) (3.0,22.0) (2.9,20.5) (2.9,19.2) (3.2,18.1) (3.5,17.1) (4.1,16.4) (4.8,16.0) (5.7,16.1) (6.9,16.7) (8.3,18.1) (9.8,20.4) (11.4,23.7) (12.8,27.8) (13.7,31.9) (13.9,35.2) (13.1,36.8) (11.6,36.6) (9.7,35.1) (7.7,32.9) (5.9,30.6) (4.5,28.3) (3.4,26.2) (2.7,24.2) (2.2,22.4) (1.9,20.7) (1.8,19.2) (1.8,17.8) (2.0,16.6) (2.2,15.5) (2.6,14.6) (3.2,13.8) (3.9,13.4) (4.8,13.3) (6.0,13.7) (7.4,14.9) (9.1,17.1) (11.1,20.6) (13.0,25.4) (14.5,30.9) (15.1,35.8) (14.5,38.6) (12.9,38.9) (10.5,37.3) (8.0,34.7) (5.8,31.9) (3.9,29.2) (2.5,26.8) (1.5,24.7) (0.8,22.8) (0.3,21.0) (-0.0,19.4) (-0.3,17.9) (-0.5,16.5) (-0.6,15.3) (-0.8,14.1) (-1.0,13.1) (-1.3,12.1) (-1.6,11.3) (-2.1,10.6) (-2.7,10.1) (-3.4,9.8) (-4.4,9.9) (-5.7,10.5) (-7.4,12.0) (-9.5,14.8) (-11.8,19.3) (-14.1,25.6) (-15.9,32.8) (-16.5,38.8) (-15.4,41.6) (-13.0,41.0) (-9.9,38.2) (-6.8,34.8) (-4.2,31.6) (-2.1,28.9) (-0.6,26.5) (0.5,24.5) (1.2,22.7) (1.9,21.1) (2.4,19.7) (2.9,18.5) (3.5,17.6) (4.1,16.9) (4.9,16.5) (5.9,16.6) (7.0,17.2) (8.4,18.6) (9.9,20.9) (11.4,24.2) (12.8,28.1) (13.6,32.1) (13.7,35.1) (12.9,36.5) (11.4,36.2) (9.5,34.7) (7.6,32.6) (5.9,30.4) (4.6,28.1) (3.5,26.0) (2.8,24.1) (2.4,22.3) (2.1,20.7) (2.1,19.2) (2.2,17.9) (2.4,16.7) (2.7,15.7) (3.2,14.9) (3.8,14.3) (4.7,14.1) (5.7,14.3) (7.1,15.2) (8.7,17.0) (10.5,20.0) (12.3,24.2) (13.8,29.3) (14.7,34.1) (14.5,37.4) (13.2,38.5) (11.2,37.4) (8.8,35.2) (6.6,32.5) (4.7,29.9) (3.2,27.5) (2.2,25.3) (1.4,23.3) (0.9,21.5) (0.6,19.9) (0.4,18.3) (0.3,16.9) (0.2,15.6) (0.2,14.4) (0.2,13.3) (0.2,12.3) (0.2,11.4) (0.3,10.5) (0.4,9.7) (0.5,8.9) (0.6,8.3) (0.8,7.7) (1.1,7.1) (1.4,6.7) (1.9,6.3) (2.6,6.1) (3.5,6.2) (4.7,6.6) (6.3,7.8) (8.3,10.2) (10.9,14.4) (13.7,21.1) (16.3,29.8) (17.8,38.4) (17.3,43.7) (14.8,44.0) (11.2,41.0) (7.3,36.9) (4.0,33.2) (1.4,30.2) (-0.5,27.8) (-1.8,25.9) (-2.8,24.3) (-3.6,23.0) (-4.3,21.9) (-4.9,21.1) (-5.6,20.7) (-6.4,20.6) (-7.3,20.9) (-8.3,21.8) (-9.3,23.2) (-10.3,25.1) (-11.1,27.4) (-11.7,29.7) (-11.9,31.6) (-11.6,32.8) (-10.9,33.2) (-9.8,32.7) (-8.7,31.6) (-7.5,30.2) (-6.5,28.6) (-5.7,27.0) (-5.2,25.5) (-4.8,24.1) (-4.7,22.8) (-4.7,21.7) (-4.9,20.8) (-5.3,20.2) (-5.9,19.8) (-6.6,19.9) (-7.5,20.4) (-8.5,21.4) (-9.6,23.1) (-10.7,25.4) (-11.6,27.9) (-12.1,30.5) (-12.2,32.5) (-11.7,33.6) (-10.8,33.7) (-9.6,32.9) (-8.3,31.6) (-7.1,30.0) (-6.1,28.3) (-5.2,26.6) (-4.7,25.0) (-4.4,23.5) (-4.2,22.2) (-4.3,21.0) (-4.6,20.1) (-5.0,19.4) (-5.6,19.0) (-6.4,19.0) (-7.3,19.5) (-8.4,20.6) (-9.6,22.4) (-10.8,24.9) (-11.8,27.7) (-12.5,30.6) (-12.6,33.0) (-12.1,34.3) (-11.1,34.4) (-9.7,33.5) (-8.3,32.1) (-6.9,30.3) (-5.8,28.4) (-4.9,26.6) (-4.3,24.8)}
\def\kviolorenza{(0.00,1.0) (0.05,1.3) (0.10,2.1) (0.15,3.7) (0.20,6.5) (0.25,11.0) (0.30,16.7) (0.35,19.6) (0.40,15.4) (0.45,7.5) (0.50,1.2) (0.55,-2.7) (0.60,-4.8) (0.65,-6.1) (0.70,-7.0) (0.75,-7.9) (0.80,-8.6) (0.85,-9.3) (0.90,-9.7) (0.95,-9.7) (1.00,-9.4) (1.05,-8.8) (1.10,-8.1) (1.15,-7.5) (1.20,-7.2) (1.25,-7.1) (1.30,-7.4) (1.35,-7.8) (1.40,-8.5) (1.45,-9.2) (1.50,-9.7) (1.55,-9.9) (1.60,-9.6) (1.65,-9.1) (1.70,-8.3) (1.75,-7.7) (1.80,-7.2) (1.85,-7.0) (1.90,-7.1) (1.95,-7.6) (2.00,-8.2) (2.05,-8.9) (2.10,-9.5) (2.15,-9.9) (2.20,-9.9) (2.25,-9.4) (2.30,-8.7) (2.35,-7.9) (2.40,-7.3) (2.45,-7.0) (2.50,-7.0) (2.55,-7.3) (2.60,-7.8) (2.65,-8.6) (2.70,-9.3) (2.75,-9.9) (2.80,-10.0) (2.85,-9.7) (2.90,-9.0) (2.95,-8.2) (3.00,-7.5) (3.05,-7.0) (3.10,-6.8) (3.15,-7.0) (3.20,-7.5) (3.25,-8.2) (3.30,-9.0) (3.35,-9.7) (3.40,-10.1) (3.45,-10.0) (3.50,-9.4) (3.55,-8.6) (3.60,-7.7) (3.65,-7.1) (3.70,-6.7) (3.75,-6.8) (3.80,-7.1) (3.85,-7.8) (3.90,-8.6) (3.95,-9.5) (4.00,-10.1) (4.05,-10.2) (4.10,-9.8) (4.15,-9.0) (4.20,-8.1) (4.25,-7.2) (4.30,-6.7) (4.35,-6.6) (4.40,-6.8) (4.45,-7.4) (4.50,-8.2) (4.55,-9.1) (4.60,-9.9) (4.65,-10.4) (4.70,-10.2) (4.75,-9.5) (4.80,-8.5) (4.85,-7.6) (4.90,-6.8) (4.95,-6.5) (5.00,-6.5) (5.05,-6.9) (5.10,-7.7) (5.15,-8.6) (5.20,-9.6) (5.25,-10.3) (5.30,-10.5) (5.35,-10.0) (5.40,-9.1) (5.45,-8.0) (5.50,-7.1) (5.55,-6.5) (5.60,-6.3) (5.65,-6.5) (5.70,-7.1) (5.75,-8.0) (5.80,-9.1) (5.85,-10.0) (5.90,-10.6) (5.95,-10.5) (6.00,-9.7) (6.05,-8.6) (6.10,-7.5) (6.15,-6.6) (6.20,-6.2) (6.25,-6.2) (6.30,-6.6) (6.35,-7.4) (6.40,-8.4) (6.45,-9.5) (6.50,-10.5) (6.55,-10.8) (6.60,-10.4) (6.65,-9.3) (6.70,-8.1) (6.75,-6.9) (6.80,-6.2) (6.85,-6.0) (6.90,-6.1) (6.95,-6.7) (7.00,-7.7) (7.05,-8.8) (7.10,-10.0) (7.15,-10.8) (7.20,-10.9) (7.25,-10.2) (7.30,-8.9) (7.35,-7.5) (7.40,-6.5) (7.45,-5.9) (7.50,-5.8) (7.55,-6.1) (7.60,-6.9) (7.65,-8.0) (7.70,-9.3) (7.75,-10.5) (7.80,-11.1) (7.85,-10.9) (7.90,-9.9) (7.95,-8.4) (8.00,-7.0) (8.05,-6.0) (8.10,-5.6) (8.15,-5.6) (8.20,-6.1) (8.25,-7.0) (8.30,-8.2) (8.35,-9.7) (8.40,-10.9) (8.45,-11.4) (8.50,-10.9) (8.55,-9.6) (8.60,-7.9) (8.65,-6.5) (8.70,-5.6) (8.75,-5.3) (8.80,-5.4) (8.85,-6.0) (8.90,-7.1) (8.95,-8.5) (9.00,-10.0) (9.05,-11.2) (9.10,-11.6) (9.15,-10.9) (9.20,-9.3) (9.25,-7.6) (9.30,-6.1) (9.35,-5.3) (9.40,-5.0) (9.45,-5.2) (9.50,-5.9) (9.55,-7.0) (9.60,-8.5) (9.65,-10.2) (9.70,-11.5) (9.75,-11.9) (9.80,-11.0) (9.85,-9.2) (9.90,-7.3) (9.95,-5.8) (10.00,-4.9) (10.05,-4.6) (10.10,-4.8) (10.15,-5.6) (10.20,-6.8) (10.25,-8.4) (10.30,-10.2) (10.35,-11.7) (10.40,-12.2) (10.45,-11.3) (10.50,-9.4) (10.55,-7.2) (10.60,-5.6) (10.65,-4.6) (10.70,-4.2) (10.75,-4.4) (10.80,-5.0) (10.85,-6.2) (10.90,-7.8) (10.95,-9.8) (11.00,-11.7) (11.05,-12.6) (11.10,-12.0) (11.15,-10.0) (11.20,-7.6) (11.25,-5.6) (11.30,-4.3) (11.35,-3.7) (11.40,-3.7) (11.45,-4.2) (11.50,-5.2) (11.55,-6.6) (11.60,-8.6) (11.65,-10.9) (11.70,-12.7) (11.75,-13.1) (11.80,-11.5) (11.85,-8.8) (11.90,-6.2) (11.95,-4.4) (12.00,-3.3) (12.05,-3.0) (12.10,-3.1) (12.15,-3.7) (12.20,-4.7) (12.25,-6.3) (12.30,-8.5) (12.35,-11.1) (12.40,-13.3) (12.45,-13.7) (12.50,-11.9) (12.55,-8.7) (12.60,-5.8) (12.65,-3.7) (12.70,-2.5) (12.75,-2.0) (12.80,-1.9) (12.85,-2.3) (12.90,-2.9) (12.95,-4.0) (13.00,-5.7) (13.05,-8.2) (13.10,-11.2) (13.15,-14.1) (13.20,-14.8) (13.25,-12.6) (13.30,-8.6) (13.35,-5.0) (13.40,-2.5) (13.45,-1.1) (13.50,-0.4) (13.55,0.0) (13.60,0.2) (13.65,0.4) (13.70,0.6) (13.75,0.9) (13.80,1.4) (13.85,2.2) (13.90,3.5) (13.95,5.6) (14.00,8.7) (14.05,12.9) (14.10,16.5) (14.15,16.4) (14.20,11.9) (14.25,6.3) (14.30,2.0) (14.35,-0.6) (14.40,-2.2) (14.45,-3.3) (14.50,-4.3) (14.55,-5.5) (14.60,-7.1) (14.65,-9.2) (14.70,-11.3) (14.75,-12.7) (14.80,-12.5) (14.85,-10.8) (14.90,-8.2) (14.95,-5.9) (15.00,-4.4) (15.05,-3.6) (15.10,-3.4) (15.15,-3.7) (15.20,-4.5) (15.25,-5.8) (15.30,-7.6) (15.35,-9.9) (15.40,-12.2) (15.45,-13.4) (15.50,-12.6) (15.55,-10.1) (15.60,-7.2) (15.65,-4.9) (15.70,-3.4) (15.75,-2.7) (15.80,-2.6) (15.85,-2.9) (15.90,-3.6) (15.95,-4.8) (16.00,-6.6) (16.05,-9.0) (16.10,-11.8) (16.15,-13.9) (16.20,-13.9) (16.25,-11.4) (16.30,-7.9) (16.35,-4.9) (16.40,-2.9) (16.45,-1.8) (16.50,-1.4) (16.55,-1.3) (16.60,-1.5) (16.65,-2.0) (16.70,-2.8) (16.75,-4.0) (16.80,-5.9) (16.85,-8.7) (16.90,-12.2) (16.95,-15.2) (17.00,-15.3) (17.05,-11.9) (17.10,-7.3) (17.15,-3.6) (17.20,-1.2) (17.25,0.2) (17.30,1.0) (17.35,1.6) (17.40,2.2) (17.45,3.1) (17.50,4.4) (17.55,6.3) (17.60,9.0) (17.65,12.2) (17.70,14.6) (17.75,14.5) (17.80,11.5) (17.85,7.5) (17.90,4.1) (17.95,2.0) (18.00,0.8) (18.05,0.2) (18.10,-0.2) (18.15,-0.4) (18.20,-0.6) (18.25,-0.9) (18.30,-1.4) (18.35,-2.1) (18.40,-3.3) (18.45,-5.2) (18.50,-8.2) (18.55,-12.1) (18.60,-15.9) (18.65,-16.6) (18.70,-12.8) (18.75,-7.3) (18.80,-2.8) (18.85,0.0) (18.90,1.7) (18.95,2.7) (19.00,3.7) (19.05,4.8) (19.10,6.4) (19.15,8.3) (19.20,10.6) (19.25,12.5) (19.30,13.1) (19.35,11.8) (19.40,9.2) (19.45,6.5) (19.50,4.6) (19.55,3.4) (19.60,3.0) (19.65,3.1) (19.70,3.6) (19.75,4.6) (19.80,6.1) (19.85,8.2) (19.90,10.8) (19.95,13.1) (20.00,13.8) (20.05,12.2) (20.10,9.1) (20.15,6.1) (20.20,3.9) (20.25,2.6) (20.30,2.0) (20.35,1.9) (20.40,2.2) (20.45,2.9) (20.50,3.9) (20.55,5.6) (20.60,7.9) (20.65,10.9) (20.70,13.8) (20.75,14.9) (20.80,12.9) (20.85,9.0) (20.90,5.3) (20.95,2.8) (21.00,1.3) (21.05,0.5) (21.10,0.1) (21.15,-0.2) (21.20,-0.3) (21.25,-0.5) (21.30,-0.8) (21.35,-1.2) (21.40,-1.8) (21.45,-2.9) (21.50,-4.6) (21.55,-7.3) (21.60,-11.2) (21.65,-15.4) (21.70,-17.2) (21.75,-14.1) (21.80,-8.4) (21.85,-3.3) (21.90,-0.1) (21.95,1.8) (22.00,3.0) (22.05,4.0) (22.10,5.2) (22.15,6.7) (22.20,8.6) (22.25,10.6) (22.30,12.3) (22.35,12.6) (22.40,11.3) (22.45,9.0) (22.50,6.7) (22.55,4.9) (22.60,3.9) (22.65,3.6) (22.70,3.8) (22.75,4.5) (22.80,5.6) (22.85,7.3) (22.90,9.5) (22.95,11.7) (23.00,13.1) (23.05,12.7) (23.10,10.6) (23.15,7.8) (23.20,5.4) (23.25,3.8) (23.30,3.1) (23.35,2.9) (23.40,3.2) (23.45,3.9) (23.50,5.1) (23.55,6.9) (23.60,9.3) (23.65,11.9) (23.70,13.7) (23.75,13.5) (23.80,11.0) (23.85,7.7) (23.90,4.9) (23.95,3.1) (24.00,2.2) (24.05,1.8) (24.10,1.9) (24.15,2.2) (24.20,3.0) (24.25,4.2) (24.30,6.0) (24.35,8.5) (24.40,11.7) (24.45,14.5) (24.50,14.9) (24.55,12.1) (24.60,8.0) (24.65,4.5) (24.70,2.1) (24.75,0.8) (24.80,0.1) (24.85,-0.3) (24.90,-0.6) (24.95,-1.0) (25.00,-1.4) (25.05,-2.1) (25.10,-3.1) (25.15,-4.8) (25.20,-7.4) (25.25,-11.0) (25.30,-14.8) (25.35,-16.5) (25.40,-13.9) (25.45,-8.8) (25.50,-4.2) (25.55,-1.1) (25.60,0.7) (25.65,1.9) (25.70,2.7) (25.75,3.7) (25.80,4.9) (25.85,6.6) (25.90,8.9) (25.95,11.4) (26.00,13.4) (26.05,13.5) (26.10,11.4) (26.15,8.3) (26.20,5.5) (26.25,3.5) (26.30,2.5) (26.35,2.1) (26.40,2.2) (26.45,2.6) (26.50,3.4) (26.55,4.7) (26.60,6.6) (26.65,9.2) (26.70,12.3) (26.75,14.5) (26.80,14.2) (26.85,11.2) (26.90,7.3) (26.95,4.2) (27.00,2.2) (27.05,1.1) (27.10,0.5) (27.15,0.3) (27.20,0.2) (27.25,0.2) (27.30,0.2) (27.35,0.3) (27.40,0.5) (27.45,0.8) (27.50,1.3) (27.55,2.1) (27.60,3.5) (27.65,5.7) (27.70,9.1) (27.75,13.7) (27.80,17.5) (27.85,16.7) (27.90,11.2) (27.95,5.1) (28.00,0.7) (28.05,-1.8) (28.10,-3.4) (28.15,-4.5) (28.20,-5.6) (28.25,-7.0) (28.30,-8.6) (28.35,-10.3) (28.40,-11.5) (28.45,-11.8) (28.50,-10.9) (28.55,-9.1) (28.60,-7.2) (28.65,-5.7) (28.70,-4.9) (28.75,-4.7) (28.80,-4.9) (28.85,-5.7) (28.90,-6.9) (28.95,-8.5) (29.00,-10.3) (29.05,-11.8) (29.10,-12.2) (29.15,-11.2) (29.20,-9.2) (29.25,-7.1) (29.30,-5.5) (29.35,-4.5) (29.40,-4.2) (29.45,-4.5) (29.50,-5.2) (29.55,-6.4) (29.60,-8.0) (29.65,-10.0) (29.70,-11.8) (29.75,-12.6) (29.80,-11.8) (29.85,-9.7) (29.90,-7.4) (29.95,-5.4) (30.00,-4.3) (30.05,-3.8) (30.10,-3.8) (30.15,-4.4) (30.20,-5.4) (30.25,-6.9) (30.30,-9.0) (30.35,-11.2) (30.40,-12.8) (30.45,-12.9) (30.50,-11.1) (30.55,-8.4) (30.60,-6.0) (30.65,-4.2) (30.70,-3.3) (30.75,-3.0) (30.80,-3.2) (30.85,-3.9) (30.90,-5.0) (30.95,-6.7) (31.00,-9.0)}
\def\kviolorenzb{(0.00,1.0) (0.05,1.3) (0.10,2.1) (0.15,3.7) (0.20,6.5) (0.25,11.0) (0.30,16.7) (0.35,19.6) (0.40,15.4) (0.45,7.5) (0.50,1.2) (0.55,-2.7) (0.60,-4.8) (0.65,-6.1) (0.70,-7.0) (0.75,-7.9) (0.80,-8.6) (0.85,-9.3) (0.90,-9.7) (0.95,-9.7) (1.00,-9.4) (1.05,-8.8) (1.10,-8.1) (1.15,-7.5) (1.20,-7.2) (1.25,-7.1) (1.30,-7.4) (1.35,-7.8) (1.40,-8.5) (1.45,-9.2) (1.50,-9.7) (1.55,-9.9) (1.60,-9.6) (1.65,-9.1) (1.70,-8.3) (1.75,-7.7) (1.80,-7.2) (1.85,-7.0) (1.90,-7.1) (1.95,-7.6) (2.00,-8.2) (2.05,-8.9) (2.10,-9.5) (2.15,-9.9) (2.20,-9.9) (2.25,-9.4) (2.30,-8.7) (2.35,-7.9) (2.40,-7.3) (2.45,-7.0) (2.50,-7.0) (2.55,-7.3) (2.60,-7.8) (2.65,-8.6) (2.70,-9.3) (2.75,-9.9) (2.80,-10.0) (2.85,-9.7) (2.90,-9.0) (2.95,-8.2) (3.00,-7.5) (3.05,-7.0) (3.10,-6.8) (3.15,-7.0) (3.20,-7.5) (3.25,-8.2) (3.30,-9.0) (3.35,-9.7) (3.40,-10.1) (3.45,-10.0) (3.50,-9.4) (3.55,-8.6) (3.60,-7.7) (3.65,-7.1) (3.70,-6.7) (3.75,-6.8) (3.80,-7.1) (3.85,-7.8) (3.90,-8.6) (3.95,-9.5) (4.00,-10.1) (4.05,-10.2) (4.10,-9.8) (4.15,-9.0) (4.20,-8.1) (4.25,-7.2) (4.30,-6.7) (4.35,-6.6) (4.40,-6.8) (4.45,-7.4) (4.50,-8.2) (4.55,-9.1) (4.60,-9.9) (4.65,-10.4) (4.70,-10.2) (4.75,-9.5) (4.80,-8.5) (4.85,-7.6) (4.90,-6.8) (4.95,-6.5) (5.00,-6.5) (5.05,-6.9) (5.10,-7.7) (5.15,-8.6) (5.20,-9.6) (5.25,-10.3) (5.30,-10.5) (5.35,-10.0) (5.40,-9.1) (5.45,-8.0) (5.50,-7.1) (5.55,-6.5) (5.60,-6.3) (5.65,-6.5) (5.70,-7.1) (5.75,-8.0) (5.80,-9.1) (5.85,-10.0) (5.90,-10.6) (5.95,-10.5) (6.00,-9.7) (6.05,-8.6) (6.10,-7.5) (6.15,-6.6) (6.20,-6.2) (6.25,-6.2) (6.30,-6.6) (6.35,-7.4) (6.40,-8.4) (6.45,-9.5) (6.50,-10.5) (6.55,-10.8) (6.60,-10.4) (6.65,-9.3) (6.70,-8.1) (6.75,-6.9) (6.80,-6.2) (6.85,-6.0) (6.90,-6.1) (6.95,-6.7) (7.00,-7.7) (7.05,-8.8) (7.10,-10.0) (7.15,-10.8) (7.20,-10.9) (7.25,-10.2) (7.30,-8.9) (7.35,-7.5) (7.40,-6.5) (7.45,-5.9) (7.50,-5.8) (7.55,-6.1) (7.60,-6.9) (7.65,-8.0) (7.70,-9.3) (7.75,-10.5) (7.80,-11.1) (7.85,-10.9) (7.90,-9.9) (7.95,-8.4) (8.00,-7.0) (8.05,-6.0) (8.10,-5.6) (8.15,-5.6) (8.20,-6.1) (8.25,-7.0) (8.30,-8.2) (8.35,-9.7) (8.40,-10.9) (8.45,-11.4) (8.50,-10.9) (8.55,-9.6) (8.60,-7.9) (8.65,-6.5) (8.70,-5.6) (8.75,-5.3) (8.80,-5.4) (8.85,-6.0) (8.90,-7.1) (8.95,-8.5) (9.00,-10.0) (9.05,-11.2) (9.10,-11.6) (9.15,-10.9) (9.20,-9.3) (9.25,-7.6) (9.30,-6.1) (9.35,-5.3) (9.40,-5.0) (9.45,-5.2) (9.50,-5.9) (9.55,-7.0) (9.60,-8.5) (9.65,-10.2) (9.70,-11.5) (9.75,-11.9) (9.80,-11.0) (9.85,-9.2) (9.90,-7.3) (9.95,-5.8) (10.00,-4.9) (10.05,-4.6) (10.10,-4.8) (10.15,-5.6) (10.20,-6.8) (10.25,-8.4) (10.30,-10.2) (10.35,-11.7) (10.40,-12.2) (10.45,-11.3) (10.50,-9.4) (10.55,-7.2) (10.60,-5.6) (10.65,-4.6) (10.70,-4.2) (10.75,-4.4) (10.80,-5.0) (10.85,-6.2) (10.90,-7.8) (10.95,-9.8) (11.00,-11.7) (11.05,-12.6) (11.10,-12.0) (11.15,-10.0) (11.20,-7.6) (11.25,-5.6) (11.30,-4.3) (11.35,-3.7) (11.40,-3.7) (11.45,-4.2) (11.50,-5.2) (11.55,-6.6) (11.60,-8.6) (11.65,-10.9) (11.70,-12.7) (11.75,-13.0) (11.80,-11.5) (11.85,-8.8) (11.90,-6.2) (11.95,-4.4) (12.00,-3.3) (12.05,-3.0) (12.10,-3.1) (12.15,-3.7) (12.20,-4.7) (12.25,-6.3) (12.30,-8.5) (12.35,-11.1) (12.40,-13.3) (12.45,-13.7) (12.50,-11.9) (12.55,-8.7) (12.60,-5.8) (12.65,-3.7) (12.70,-2.5) (12.75,-2.0) (12.80,-1.9) (12.85,-2.3) (12.90,-2.9) (12.95,-4.0) (13.00,-5.7) (13.05,-8.2) (13.10,-11.2) (13.15,-14.1) (13.20,-14.8) (13.25,-12.6) (13.30,-8.6) (13.35,-5.0) (13.40,-2.5) (13.45,-1.1) (13.50,-0.4) (13.55,0.0) (13.60,0.2) (13.65,0.4) (13.70,0.6) (13.75,0.9) (13.80,1.4) (13.85,2.2) (13.90,3.5) (13.95,5.6) (14.00,8.7) (14.05,12.9) (14.10,16.5) (14.15,16.4) (14.20,11.9) (14.25,6.2) (14.30,2.0) (14.35,-0.6) (14.40,-2.2) (14.45,-3.3) (14.50,-4.3) (14.55,-5.5) (14.60,-7.2) (14.65,-9.2) (14.70,-11.3) (14.75,-12.7) (14.80,-12.5) (14.85,-10.7) (14.90,-8.2) (14.95,-5.9) (15.00,-4.4) (15.05,-3.6) (15.10,-3.4) (15.15,-3.7) (15.20,-4.5) (15.25,-5.8) (15.30,-7.6) (15.35,-9.9) (15.40,-12.2) (15.45,-13.4) (15.50,-12.6) (15.55,-10.1) (15.60,-7.2) (15.65,-4.9) (15.70,-3.4) (15.75,-2.7) (15.80,-2.6) (15.85,-2.9) (15.90,-3.6) (15.95,-4.8) (16.00,-6.6) (16.05,-9.0) (16.10,-11.8) (16.15,-13.9) (16.20,-13.9) (16.25,-11.4) (16.30,-7.9) (16.35,-4.9) (16.40,-2.9) (16.45,-1.8) (16.50,-1.4) (16.55,-1.3) (16.60,-1.5) (16.65,-1.9) (16.70,-2.7) (16.75,-4.0) (16.80,-5.9) (16.85,-8.7) (16.90,-12.2) (16.95,-15.1) (17.00,-15.3) (17.05,-12.1) (17.10,-7.4) (17.15,-3.7) (17.20,-1.2) (17.25,0.1) (17.30,1.0) (17.35,1.6) (17.40,2.2) (17.45,3.1) (17.50,4.4) (17.55,6.4) (17.60,9.0) (17.65,12.2) (17.70,14.6) (17.75,14.5) (17.80,11.5) (17.85,7.5) (17.90,4.1) (17.95,2.0) (18.00,0.8) (18.05,0.2) (18.10,-0.1) (18.15,-0.3) (18.20,-0.5) (18.25,-0.8) (18.30,-1.2) (18.35,-1.9) (18.40,-2.9) (18.45,-4.6) (18.50,-7.2) (18.55,-11.0) (18.60,-15.1) (18.65,-17.0) (18.70,-14.2) (18.75,-8.7) (18.80,-3.7) (18.85,-0.5) (18.90,1.5) (18.95,2.7) (19.00,3.6) (19.05,4.7) (19.10,6.2) (19.15,8.0) (19.20,10.2) (19.25,12.2) (19.30,13.0) (19.35,12.0) (19.40,9.6) (19.45,7.0) (19.50,5.0) (19.55,3.8) (19.60,3.2) (19.65,3.3) (19.70,3.8) (19.75,4.7) (19.80,6.2) (19.85,8.3) (19.90,10.7) (19.95,12.9) (20.00,13.6) (20.05,12.1) (20.10,9.1) (20.15,6.2) (20.20,4.1) (20.25,2.9) (20.30,2.3) (20.35,2.3) (20.40,2.7) (20.45,3.5) (20.50,4.7) (20.55,6.6) (20.60,9.1) (20.65,12.0) (20.70,14.2) (20.75,14.1) (20.80,11.3) (20.85,7.6) (20.90,4.5) (20.95,2.5) (21.00,1.4) (21.05,0.9) (21.10,0.7) (21.15,0.7) (21.20,0.9) (21.25,1.3) (21.30,1.9) (21.35,2.9) (21.40,4.5) (21.45,6.9) (21.50,10.4) (21.55,14.4) (21.60,16.6) (21.65,14.6) (21.70,9.6) (21.75,4.7) (21.80,1.3) (21.85,-0.7) (21.90,-1.9) (21.95,-2.8) (22.00,-3.8) (22.05,-5.0) (22.10,-6.7) (22.15,-8.9) (22.20,-11.4) (22.25,-13.2) (22.30,-13.3) (22.35,-11.3) (22.40,-8.3) (22.45,-5.6) (22.50,-3.8) (22.55,-2.8) (22.60,-2.4) (22.65,-2.5) (22.70,-3.0) (22.75,-4.0) (22.80,-5.4) (22.85,-7.5) (22.90,-10.3) (22.95,-13.0) (23.00,-14.4) (23.05,-13.1) (23.10,-9.8) (23.15,-6.3) (23.20,-3.7) (23.25,-2.1) (23.30,-1.3) (23.35,-1.0) (23.40,-0.9) (23.45,-1.1) (23.50,-1.5) (23.55,-2.2) (23.60,-3.3) (23.65,-4.9) (23.70,-7.5) (23.75,-11.0) (23.80,-14.7) (23.85,-16.2) (23.90,-13.7) (23.95,-8.9) (24.00,-4.4) (24.05,-1.4) (24.10,0.4) (24.15,1.5) (24.20,2.3) (24.25,3.1) (24.30,4.2) (24.35,5.8) (24.40,7.9) (24.45,10.6) (24.50,13.1) (24.55,14.1) (24.60,12.6) (24.65,9.4) (24.70,6.1) (24.75,3.7) (24.80,2.3) (24.85,1.6) (24.90,1.4) (24.95,1.6) (25.00,2.0) (25.05,2.7) (25.10,3.9) (25.15,5.7) (25.20,8.3) (25.25,11.7) (25.30,14.7) (25.35,15.3) (25.40,12.5) (25.45,8.0) (25.50,4.2) (25.55,1.7) (25.60,0.3) (25.65,-0.6) (25.70,-1.1) (25.75,-1.6) (25.80,-2.3) (25.85,-3.3) (25.90,-4.8) (25.95,-7.1) (26.00,-10.1) (26.05,-13.4) (26.10,-15.3) (26.15,-14.0) (26.20,-10.1) (26.25,-5.9) (26.30,-2.8) (26.35,-1.0) (26.40,0.0) (26.45,0.6) (26.50,1.1) (26.55,1.6) (26.60,2.3) (26.65,3.3) (26.70,4.9) (26.75,7.3) (26.80,10.5) (26.85,13.9) (26.90,15.7) (26.95,13.9) (27.00,9.6) (27.05,5.3) (27.10,2.3) (27.15,0.5) (27.20,-0.6) (27.25,-1.2) (27.30,-1.8) (27.35,-2.5) (27.40,-3.6) (27.45,-5.1) (27.50,-7.4) (27.55,-10.3) (27.60,-13.4) (27.65,-15.0) (27.70,-13.6) (27.75,-9.8) (27.80,-5.9) (27.85,-3.0) (27.90,-1.3) (27.95,-0.4) (28.00,0.1) (28.05,0.4) (28.10,0.7) (28.15,1.0) (28.20,1.6) (28.25,2.3) (28.30,3.6) (28.35,5.5) (28.40,8.5) (28.45,12.4) (28.50,15.9) (28.55,16.2) (28.60,12.3) (28.65,7.0) (28.70,2.8) (28.75,0.2) (28.80,-1.3) (28.85,-2.3) (28.90,-3.2) (28.95,-4.3) (29.00,-5.7) (29.05,-7.7) (29.10,-10.1) (29.15,-12.4) (29.20,-13.6) (29.25,-12.6) (29.30,-9.9) (29.35,-6.9) (29.40,-4.6) (29.45,-3.2) (29.50,-2.5) (29.55,-2.4) (29.60,-2.6) (29.65,-3.3) (29.70,-4.5) (29.75,-6.2) (29.80,-8.6) (29.85,-11.5) (29.90,-13.9) (29.95,-14.3) (30.00,-11.9) (30.05,-8.3) (30.10,-5.0) (30.15,-2.8) (30.20,-1.6) (30.25,-1.0) (30.30,-0.8) (30.35,-0.9) (30.40,-1.1) (30.45,-1.5) (30.50,-2.2) (30.55,-3.4) (30.60,-5.1) (30.65,-7.8) (30.70,-11.5) (30.75,-15.1) (30.80,-16.3) (30.85,-13.3) (30.90,-8.2) (30.95,-3.8) (31.00,-0.9)}
\def\kviorctwo{(0.0,0.00) (0.2,0.05) (0.4,0.16) (0.6,0.34) (0.8,0.55) (1.0,0.79) (1.2,1.04) (1.4,1.31) (1.6,1.58) (1.8,1.86) (2.0,2.13) (2.2,2.40) (2.4,2.67) (2.6,2.93) (2.8,3.18) (3.0,3.43) (3.2,3.67) (3.4,3.90) (3.6,4.13) (3.8,4.35) (4.0,4.56) (4.2,4.76) (4.4,4.95) (4.6,5.14) (4.8,5.32) (5.0,5.50) (5.2,5.66) (5.4,5.83) (5.6,5.98) (5.8,6.13) (6.0,6.28) (6.2,6.42) (6.4,6.55) (6.6,6.68) (6.8,6.81) (7.0,6.92) (7.2,7.04) (7.4,7.15) (7.6,7.26) (7.8,7.36) (8.0,7.46) (8.2,7.55) (8.4,7.65) (8.6,7.73) (8.8,7.82) (9.0,7.90) (9.2,7.98) (9.4,8.06) (9.6,8.13) (9.8,8.20) (10.0,8.27) (10.2,8.33) (10.4,8.39) (10.6,8.45) (10.8,8.51) (11.0,8.57) (11.2,8.62) (11.4,8.67) (11.6,8.72) (11.8,8.77) (12.0,8.82) (12.2,8.86) (12.4,8.90) (12.6,8.94) (12.8,8.98) (13.0,9.02) (13.2,9.06) (13.4,9.09) (13.6,9.13) (13.8,9.16) (14.0,9.19) (14.2,9.22) (14.4,9.25) (14.6,9.28) (14.8,9.31) (15.0,9.33) (15.2,9.36) (15.4,9.38) (15.6,9.40) (15.8,9.43) (16.0,9.45) (16.2,9.47)}

\bkchapter{6}{14}{微分方程式がいくつも：\\連立微分方程式}{\bklevel{5}}{微分方程式がいくつも：連立微分方程式}
\bkmichishirube{3-6,4-9,6-2,6-4,6-6,6-12}{7-8,7-9,7-10,9-10,9-11}

\begin{goal}
\begin{itemize}[label={\color{\bkTT}●}, leftmargin=1.6\zw, itemsep=0pt, topsep=1pt]
\item \bklevel{3}\enspace からみあう2つの量の変化を、表計算で2つの列を同時に進めて計算できる。
\item \bklevel{4}\enspace 和と差（重みをつけた和と差）に分けて、2つの式を 6-4 の形にして解ける。
\item \bklevel{5}\enspace 振動モードや消去法で、連立微分方程式の答えを式で求められる。
\end{itemize}
\end{goal}

\begin{hatena}
\noindent\begin{minipage}[c]{0.56\linewidth}
\raggedright
同じ形の2つの水そうを、底の近くで細い管でつなぎます。はじめ、左の水の深さは\qty{20}{cm}、右は空っぽです。

管をあけると、水は左から右へ流れます。左の水位は下がり、右の水位は上がります。2つの水位は、それぞれどう変わるでしょう。最後は何cmになるでしょう。
\end{minipage}\hfill
\begin{minipage}[c]{0.4\linewidth}
\centering
\begin{tikzpicture}[x=0.78mm, y=0.78mm]
  \kviotank{(0,0)}{16}{26}{20}
  \kviotank{(30,0)}{16}{26}{0}
  \fill[bkBlue!22] (16,1.4) rectangle (30,3.2);
  \draw[bkInk!75, line width=0.7pt] (16,1.4) -- (30,1.4) (16,3.2) -- (30,3.2);
  \draw[bkBlue!80!black, line width=1.1pt, -{Stealth[length=3.5pt]}] (18.5,6.5) -- (27.5,6.5);
  \draw[bkInk, line width=0.6pt, {Stealth[length=3pt]}-{Stealth[length=3pt]}] (-3,0) -- (-3,20);
  \node[font=\sffamily\footnotesize, anchor=east] at (-4,10) {20\,cm};
  \node[font=\sffamily\footnotesize, anchor=north] at (8,-1.5) {左};
  \node[font=\sffamily\footnotesize, anchor=north] at (38,-1.5) {右};
  \node[font=\sffamily\footnotesize, anchor=south] at (23,8) {水};
\end{tikzpicture}
\end{minipage}
\end{hatena}

\noindent これまでの章では、変わる量は1つでした。お湯の温度 $T$、ばねにつけたおもりの位置 $x$、コンデンサーの電圧 $V$ です。ところが、身のまわりには、2つ以上の量が、おたがいに相手の変わり方を決めあうことがたくさんあります。つながった水そう、となり合う2つの部屋の温度、ばねでつながった2両の車両、親から子へ変わる放射性物質などです。

この章では、そういう「微分方程式の組」を、まず表計算で解き、つぎに式で解きます。★★★★★ の章ですが、「目でわかる」「手でわかる」の段は、中学生でも読めるように書きました。式の段は、1回で分からなくても大丈夫です。表計算の段までで、未来はもう計算できます。

%% ---------------- やってみよう・思い出そう ----------------
\begin{yattemiyou}[つないだ水そうの水位をはかろう]
\kviobox{用意するもの}
同じ形の円柱の入れ物2つ（底面積が\qty{100}{cm^2}ほど）、細いチューブ（長さ\qty{1}{m}ほど）、せんたくばさみ、ものさし、タイマー
\begin{enumerate}[label={\bkcircled[bkGreen]{\arabic*}}, leftmargin=2.4\zw, itemsep=1pt, topsep=3pt]
\item 2つの入れ物の底の近くに穴をあけ、チューブでつなぐ。チューブはせんたくばさみでとめておく。
\item 左の入れ物に、水を深さ\qty{20}{cm}まで入れる。右は空のまま。
\item せんたくばさみを外し、1分ごとに2つの水位をはかって表に書く。
\item 2つの水位の「合計」と「差」も計算しよう。何か気づくかな？
\end{enumerate}
\begin{bkcenter}
\begin{bktable}{colspec={Q[l,m,wd=27mm] *{8}{Q[c,m,wd=8.6mm]}}, rowsep=1pt, colsep=2pt}
時刻 $t$（分） & 0 & 1 & 2 & 3 & 4 & 5 & 6 & 8 \\
左の水位（cm） & 20.0 & 18.2 & 16.7 & 15.5 & 14.5 & 13.7 & 13.0 & 12.0 \\
右の水位（cm） & 0.0 & 1.8 & 3.3 & 4.5 & 5.5 & 6.3 & 7.0 & 8.0 \\
合計（cm） & 20.0 & 20.0 & 20.0 & 20.0 & 20.0 & 20.0 & 20.0 & 20.0 \\
差（cm） & 20.0 & 16.4 & 13.4 & 11.0 & 9.0 & 7.4 & 6.0 & 4.0 \\
\end{bktable}
\bktabcaption{ソラさんがはかった結果}
\end{bkcenter}
\end{yattemiyou}

\noindent 表から、2つのことに気づきます。
\begin{enumerate}[label={\bkcircled[\bkTT]{\arabic*}}, leftmargin=2.4\zw, itemsep=2pt, topsep=2pt]
\item \textbf{合計は、いつも\qty{20}{cm}。}\par\noindent 水は管を通って移るだけで、ふえもへりもしないからです（2つの入れ物は同じ形なので、水位の合計が水の量を表します）。
\item \textbf{差は、1分ごとに、ほぼ同じ割合でへる。}\par\noindent $16.4\div20.0=0.82$、\ $13.4\div16.4\approx0.82$、\ $11.0\div13.4\approx0.82$ です。毎分、前の約0.82倍になっています。3-11 の「同じ時間ごとに同じ割合でへる」、つまり 6-4 の指数関数のへり方です。
\end{enumerate}
左の水位だけ、右の水位だけを見ると、へり方もふえ方もはっきりしません。ところが「合計」と「差」に組みかえると、とても簡単なきまりが見えてきました。この組みかえが、この章のいちばん大切な考えです。

\begin{bkfigure}
\begin{tikzpicture}[x=1mm, y=0.72mm]
  \foreach \t/\a/\b/\xs in {0/20/0/0, 2/16.7/3.3/44, 8/12.0/8.0/88} {
    \begin{scope}[xshift=\xs mm]
      \kviotank{(0,0)}{12}{24}{\a}
      \kviotank{(18,0)}{12}{24}{\b}
      \draw[bkInk!75, line width=0.6pt] (12,1) -- (18,1);
      \node[font=\sffamily\footnotesize, anchor=north] at (15,-1.5) {$t=\t$ 分};
    \end{scope}}
  \draw[bkGreen!70!black, line width=0.9pt, dash pattern=on 2pt off 1.5pt] (-2,10) -- (124,10);
  \node[font=\sffamily\footnotesize, text=bkGreen!60!black, anchor=south east] at (124,25) {合計の半分（10\,cm）};
\end{tikzpicture}
\bkfigcaption{0分、2分、8分の水位。合計の半分の10\,cm（点線）から見ると、左は上に、右は下に、いつも同じだけはなれている。そのはなれ（差の半分）が、だんだん小さくなる。}
\end{bkfigure}

\begin{omoidasou}
\bkrecall{3-6}{連立方程式：2つの式から、代入法（1つの式を「$y=\cdots$」の形にして、もう1つに入れる）や加減法（式どうしをたす・ひく）で、文字を1つずつ消して解く。}
\bkrecall{6-2}{オイラー法：次の値＝今の値＋（今の変わる速さ）×$\Delta t$。2回微分の式は、位置 $x$ と速度 $v$ の2つの列で進める（速度を先に新しくし、その新しい速度で位置を進める）。}
\bkrecall{6-4}{$\dfrac{\dd y}{\dd t}=-\dfrac{y}{\tau}$ の答えは $y=y_0e^{-t/\tau}$。$\tau$ は時定数。半分になるまでの時間は $\tau\ln2$。}
\bkrecall{6-6}{$m\ddot x=-kx$ の答えは $x=A\cos\omega_{\mathrm n}t+B\sin\omega_{\mathrm n}t$。固有角振動数は $\omega_{\mathrm n}=\sqrt{k/m}$。}
\bkrecall{6-12}{答えを $e^{st}$ と置いて、特性方程式を作る。解が2つ（$s_1$、$s_2$）なら $Ae^{s_1t}+Be^{s_2t}$。重解 $s$ なら $(A+Bt)e^{st}$。}
\bkrecall{4-9}{振動数が少しちがう2つの波をたすと、大きさがゆっくり大きくなったり小さくなったりする（うなり）。}
\end{omoidasou}

%% ---------------- 数学の道具 ----------------
\bksection{doguu}{数学の道具}

\bkstage{me}{目でわかる}{からみあう2つの量}{\bklevel{3}}

\bksubsection{1}{つながった2つの水そう}

\label{bk:6-14:suiso}やってみようの水そうで、何が水の流れを決めているかを考えます。管の中の水は、水位の高いほうから低いほうへ流れます。水位の差が大きいほど強く押されて、たくさん流れます。細い管の中をゆっくり流れるときは、1分に流れる水の量は、水位の差に比例します。これは実験で確かめられたきまりです（物理シリーズ『力学』で、管の中の流れとしてくわしく学びます）。

ここでは、水位の差が\qty{1}{cm}あたり、1分に\qty{10}{cm^3}流れる管だとします。左の水位を $h_1$〔cm〕、右の水位を $h_2$〔cm〕とすると、1分に流れる水の量は $10(h_1-h_2)$〔cm³〕です。

\begin{bkfigure}
\begin{tikzpicture}[x=1mm, y=1mm]
  \kviotank{(0,0)}{24}{34}{26}
  \kviotank{(52,0)}{24}{34}{8}
  \fill[bkBlue!22] (24,1.4) rectangle (52,3.4);
  \draw[bkInk!75, line width=0.7pt] (24,1.4) -- (52,1.4) (24,3.4) -- (52,3.4);
  \draw[bkBlue!80!black, line width=1.3pt, -{Stealth[length=4pt]}] (29,8) -- (47,8);
  \node[font=\sffamily\footnotesize, anchor=south, align=center] at (38,10) {1分に\\$10(h_1-h_2)$\,cm³};
  \draw[bkInk, line width=0.6pt, {Stealth[length=3pt]}-{Stealth[length=3pt]}] (-3,0) -- (-3,26);
  \node[font=\sffamily\footnotesize, anchor=east] at (-4.5,13) {$h_1$};
  \draw[bkInk, line width=0.6pt, {Stealth[length=3pt]}-{Stealth[length=3pt]}] (79,0) -- (79,8);
  \node[font=\sffamily\footnotesize, anchor=west] at (80.5,4) {$h_2$};
  \draw[bkRed, line width=1.1pt, -{Stealth[length=4pt]}] (12,31) -- (12,24);
  \node[font=\sffamily\footnotesize, anchor=south, text=bkRed!80!black] at (12,31.5) {下がる};
  \draw[bkRed, line width=1.1pt, -{Stealth[length=4pt]}] (64,12) -- (64,19);
  \node[font=\sffamily\footnotesize, anchor=south, text=bkRed!80!black] at (64,19.5) {上がる};
  \node[font=\sffamily\footnotesize, anchor=north, align=center] at (12,-1.5) {左の水そう\\底面積 \qty{100}{cm^2}};
  \node[font=\sffamily\footnotesize, anchor=north, align=center] at (64,-1.5) {右の水そう\\底面積 \qty{100}{cm^2}};
\end{tikzpicture}
\bkfigcaption{左から右へ流れる水の量は、水位の差 $h_1-h_2$ に比例する。左で出ていった水が、そのまま右に入る。}
\end{bkfigure}

左の水そうからは、1分に $10(h_1-h_2)$\,cm³ の水が出ていきます。底面積が\qty{100}{cm^2}なので、水位は1分に
\kvioline{$10(h_1-h_2)\div100=0.1(h_1-h_2)$\quad 〔cm〕}
下がります。「水の体積÷底面積＝高さ」だからです。右の水そうには、同じ量の水が入るので、同じだけ上がります。これを、変わる速さの式に書きます。
\kvioframe[110mm]{$\dfrac{\dd h_1}{\dd t}=-0.1(h_1-h_2)$\qquad $\dfrac{\dd h_2}{\dd t}=0.1(h_1-h_2)$}
読み方は「左の水位の変わる速さ（cm/分）は、水位の差の0.1倍で、下がる向き」「右の水位の変わる速さは、水位の差の0.1倍で、上がる向き」です。
\begin{itemize}[leftmargin=1.2\zw, itemsep=2pt, topsep=2pt, label={\color{\bkTT}●}]
\item \textbf{0.1 はどこから？}\par\noindent 管が流す量（差\qty{1}{cm}あたり毎分\qty{10}{cm^3}）を、底面積\qty{100}{cm^2}でわった数です。単位は「1/分」です。
\item \textbf{マイナスはなぜ？}\par\noindent $h_1>h_2$ のとき、左の水位は下がるので、変わる速さは負です。右は上がるので正です。
\end{itemize}

ここで大切なことがあります。左の水位の式の右辺に、\textbf{右の水位 $h_2$ が入っています}。右の水位の式には、左の水位 $h_1$ が入っています。だから、1つの式だけを取り出して解くことはできません。6-4 の冷めるお湯では、右辺には温度 $T$ だけが入っていたので、その式だけで解けました。

\begin{bkfigure}
\begin{tikzpicture}[x=1mm, y=1mm]
  \node[bk node, minimum width=26mm] (a) at (0,0) {左の水位 $h_1$};
  \node[bk node, minimum width=26mm] (b) at (54,0) {右の水位 $h_2$};
  \draw[bk link, bkBlue, bend left=20] (a.north east) to node[bk link label, above] {$h_1$ が $h_2$ の変わり方を決める} (b.north west);
  \draw[bk link, bkOrange, bend left=20] (b.south west) to node[bk link label, below] {$h_2$ が $h_1$ の変わり方を決める} (a.south east);
  \node[bk node, minimum width=26mm] (c) at (100,0) {お湯の温度 $T$};
  \draw[bk link, bkGray] ([xshift=4mm]c.north) .. controls +(6,13) and +(-6,13) .. node[bk link label, above] {自分で決める} ([xshift=-4mm]c.north);
  \node[font=\sffamily\footnotesize, text=bkGray, anchor=north] at (100,-6.5) {6-4（式が1つ）};
\end{tikzpicture}
\bkfigcaption{左：2つの水位は、おたがいに相手の変わり方を決めあう。右：冷めるお湯は、自分の温度だけで変わり方が決まる。}
\end{bkfigure}

このように、いくつかの量の変わる速さが、おたがいの量で決まるような微分方程式の組を、\textbf{連立微分方程式}（れんりつびぶんほうていしき）といいます。\bkidx{れんりつびぶんほうていしき}{連立微分方程式}3-6 の連立方程式は「2つの式を同時にみたす数の組」をさがしました。連立微分方程式は「2つの式を同時にみたす\textbf{関数の組}」をさがします。

\kviohead{最後はどうなる？}
水がいつまでも流れ続けることはありません。流れが止まるのは、2つの水位が同じになって、差が0になったときです。このとき、どちらの式の右辺も0になり、水位は変わらなくなります（6-5 の終端速度で、変わる速さが0になる所をさがしたのと同じ考えです）。水の合計は\qty{20}{cm}のままなので、
\kvioline{最後の水位＝$20\div2=10$\quad 〔cm〕}
です。では、どんな速さで10\,cm に近づくのでしょう。それは、③と⑤で、表計算と式で求めます。

\bksubsection{2}{ばねでつながった2つのおもり}

\label{bk:6-14:omori}こんどは、動くものです。なめらかな床の上に、同じ質量 $m$ の2つの台車を置き、図のように3本のばねでつなぎます。両はしのばねのばね定数は $k$、まん中の「つなぎのばね」のばね定数は $K$ です。

\begin{bkfigure}
\begin{tikzpicture}[x=1mm, y=1mm]
  \draw[bkInk, line width=0.9pt] (0,0) -- (124,0);
  \kviowall{(2,0)}{14}{1}
  \kviowall{(122,0)}{14}{-1}
  \kviospring{(2,7)}{(30,7)}
  \kviocart{(30,0)}{$m$}{bkBlue}
  \kviospring{(42,7)}{(80,7)}
  \kviocart{(80,0)}{$m$}{bkOrange}
  \kviospring{(92,7)}{(122,7)}
  \node[font=\sffamily\footnotesize, anchor=south] at (16,10) {ばね $k$};
  \node[font=\sffamily\footnotesize, anchor=south] at (61,10) {つなぎのばね $K$};
  \node[font=\sffamily\footnotesize, anchor=south] at (107,10) {ばね $k$};
  \draw[bkInk, line width=0.5pt, dash pattern=on 1.5pt off 1.5pt] (36,-2) -- (36,-9) (86,-2) -- (86,-9);
  \draw[bkBlue!80!black, line width=1.1pt, -{Stealth[length=4pt]}] (36,-6) -- (46,-6);
  \draw[bkOrange!80!black, line width=1.1pt, -{Stealth[length=4pt]}] (86,-6) -- (96,-6);
  \node[font=\sffamily\footnotesize, anchor=west, text=bkBlue!80!black] at (47,-6) {$x_1$};
  \node[font=\sffamily\footnotesize, anchor=west, text=bkOrange!80!black] at (97,-6) {$x_2$};
  \node[font=\sffamily\footnotesize, anchor=north] at (36,-10) {台車1のつり合いの位置};
  \node[font=\sffamily\footnotesize, anchor=north] at (86,-10) {台車2のつり合いの位置};
\end{tikzpicture}
\bkfigcaption{2つの台車と3本のばね。$x_1$、$x_2$ は、それぞれの台車の、つり合いの位置からのずれ（右向きを正）。}
\end{bkfigure}

まず、どんな動きになるかを見てみます。両はしのばねを $k=\qty{100}{N/m}$、まん中のばねを弱く $K=\qty{10}{N/m}$、台車をそれぞれ $m=\qty{1}{kg}$ とします。台車2を止めておいて、台車1だけを右へ\qty{10}{cm}引き、2つとも静かに手をはなします。

\begin{bkfigure}
\begin{tikzpicture}
\begin{axis}[bkplot, axis x line=bottom, axis y line=left, width=124mm, height=34mm, xmin=0, xmax=7.2, ymin=-0.13, ymax=0.13, xtick={0,1,2,3,4,5,6,7}, ytick={-0.1,0,0.1}, xlabel={$t$（s）}, ylabel={$x_1$（m）}, yticklabel style={/pgf/number format/fixed}]
  \addplot[bkthin, line width=0.6pt, domain=0:7.1] {0};
  \addplot[bkthin, dashed, domain=0:7, samples=80] {0.1*cos(deg(0.4772*x))};
  \addplot[bkthin, dashed, domain=0:7, samples=80] {-0.1*cos(deg(0.4772*x))};
  \addplot[bkcurve, line width=1.1pt, domain=0:7, samples=900] {0.05*(cos(deg(10*x))+cos(deg(10.9545*x)))};
\end{axis}
\end{tikzpicture}\par\vspace{2mm}
\begin{tikzpicture}
\begin{axis}[bkplot, axis x line=bottom, axis y line=left, width=124mm, height=34mm, xmin=0, xmax=7.2, ymin=-0.13, ymax=0.13, xtick={0,1,2,3,4,5,6,7}, ytick={-0.1,0,0.1}, xlabel={$t$（s）}, ylabel={$x_2$（m）}, yticklabel style={/pgf/number format/fixed}]
  \addplot[bkthin, line width=0.6pt, domain=0:7.1] {0};
  \addplot[bkthin, dashed, domain=0:7, samples=80] {0.1*sin(deg(0.4772*x))};
  \addplot[bkthin, dashed, domain=0:7, samples=80] {-0.1*sin(deg(0.4772*x))};
  \addplot[bkcurve2, line width=1.1pt, domain=0:7, samples=900] {0.05*(cos(deg(10*x))-cos(deg(10.9545*x)))};
\end{axis}
\end{tikzpicture}
\bkfigcaption{上が台車1、下が台車2の位置（表計算で $\Delta t=\qty{0.001}{s}$ として計算したもの）。ゆれが台車1から台車2へうつり、約3.3秒で台車1はほとんど止まる。約6.6秒で、ゆれはまた台車1にもどる。}
\label{bk:6-14:fig-uneri}
\end{bkfigure}

図\ref{bk:6-14:fig-uneri}を見ると、はじめは台車1だけが大きくゆれ、台車2はほとんど動きません。ところが、台車1のゆれはだんだん小さくなり、そのかわりに台車2のゆれが大きくなります。約3.3秒で、ゆれはすっかり台車2にうつり、そのあとまた台車1にもどってきます。弱いばね1本で、ゆれが行ったり来たりするのです。点線のつつみの形は、4-9 のうなりとそっくりです。なぜこうなるかは、⑥で式から説明します。

\kviohead{2つの台車の運動方程式}
台車1と台車2にはたらく力を、1つずつ書きます。ばねの力は「のび（またはちぢみ）×ばね定数」で、のびをもどす向きにはたらきます（6-6）。
\begin{enumerate}[label={\bkcircled[\bkTT]{\arabic*}}, leftmargin=2.4\zw, itemsep=2pt, topsep=2pt]
\item \textbf{左のばね}\quad 台車1が右へ $x_1$ ずれると、左のばねは $x_1$ だけのびて、台車1を左へ引きます。力は $-kx_1$。
\item \textbf{つなぎのばね}\quad のびは $x_2-x_1$ です。たとえば $x_1=\qty{2}{cm}$、$x_2=\qty{5}{cm}$ なら、右はしが左はしより\qty{3}{cm}多く右へ動いたので、のびは\qty{3}{cm}。のびたばねは両はしを内側へ引くので、台車1は右へ $+K(x_2-x_1)$、台車2は左へ $-K(x_2-x_1)$ の力を受けます。
\item \textbf{右のばね}\quad 台車2が右へ $x_2$ ずれると、右のばねは $x_2$ だけちぢんで、台車2を左へ押します。力は $-kx_2$。
\end{enumerate}
2つが同じだけ動いたとき（$x_1=x_2$）は、つなぎのばねののびは0です。$x_2-x_1$ という書き方は、このことも正しく表しています。力の合計＝質量×加速度（5-5）なので、
\kvioframe[118mm]{$m\ddot x_1=-kx_1+K(x_2-x_1)$\qquad $m\ddot x_2=-kx_2-K(x_2-x_1)$}
ここでも、台車1の式に $x_2$、台車2の式に $x_1$ が入っています。2つのおもりの動きは、連立微分方程式で決まります。

\bkstage{te}{手でわかる}{表計算で、2つを同時に進める}{\bklevel{3}}

\bksubsection{3}{2つの列を、同じ行の値で進める}

\label{bk:6-14:hyou}連立微分方程式も、6-2 のオイラー法で解けます。列を2つにして、2つの水位を同時に進めるだけです。
\kvioline{次の $h_1$＝今の $h_1$＋（$-0.1\times$ 今の差）$\times\Delta t$}
\kvioline{次の $h_2$＝今の $h_2$＋（$\phantom{-}0.1\times$ 今の差）$\times\Delta t$}
大切なのは、\textbf{2つとも「今の行」の値で計算する}ことです。左の水位を新しくしてから、その新しい値で右を計算してはいけません（まちがえやすいところ①）。表計算では、次のように入れます（F2 に $\Delta t$ を入れる）。
\begin{itemize}[leftmargin=1.2\zw, itemsep=1pt, topsep=2pt, label={\color{\bkTT}●}]
\item D2 に \texttt{=B2-C2}（差）、E2 に \texttt{=B2+C2}（合計）
\item A3 に \texttt{=A2+\$F\$2}
\item B3 に \texttt{=B2-0.1*D2*\$F\$2}、C3 に \texttt{=C2+0.1*D2*\$F\$2}
\item 3行めの A〜E を、下へコピーする
\end{itemize}
B3 と C3 は、どちらも2行めの差 D2 を使っています。だから、2つとも「今の行」の値で計算されます。

\begin{bkcenter}
\bksheet{5}{17}{%
1 & $t$（分） & $h_1$ & $h_2$ & 差 & 合計 \\
2 & 0 & 20 & 0 & 20 & 20 \\
3 & 1 & 18 & 2 & 16 & 20 \\
4 & 2 & 16.4 & 3.6 & 12.8 & 20 \\
5 & 3 & 15.12 & 4.88 & 10.24 & 20 \\
6 & 4 & 14.096 & 5.904 & 8.192 & 20 \\
7 & 5 & 13.277 & 6.723 & 6.554 & 20 \\
12 & 10 & 11.074 & 8.926 & 2.147 & 20 \\}
\bktabcaption{$\Delta t=1$ 分で進めた表（8〜11行めは省略）。合計はぴったり20のまま。}
\end{bkcenter}

2行めから3行めへは、差が20なので、左は $0.1\times20\times1=2$ 下がって18、右は2上がって2になります。3行めから4行めへは、差が16なので、1.6ずつ動いて16.4と3.6です。

\begin{bkfigure}
\begin{tikzpicture}
\begin{axis}[bkplot, width=104mm, height=52mm, xmin=0, xmax=10.8, ymin=0, ymax=22.5, xtick={2,4,6,8,10}, ytick={5,10,15,20}, xlabel={$t$（分）}, ylabel={水位（cm）}]
  \addplot[bkthin, dashed, domain=0:10.5] {10};
  \addplot[bkcurve, domain=0:10.5, samples=60] {10+10*exp(-0.2*x)};
  \addplot[bkcurve2, domain=0:10.5, samples=60] {10-10*exp(-0.2*x)};
  \addplot[only marks, mark=*, mark size=1.7pt, bkBlue!60!black] coordinates {(0,20)(1,18)(2,16.4)(3,15.12)(4,14.096)(5,13.277)(6,12.621)(7,12.097)(8,11.678)(9,11.342)(10,11.074)};
  \addplot[only marks, mark=square*, mark size=1.5pt, bkOrange!70!black] coordinates {(0,0)(1,2)(2,3.6)(3,4.88)(4,5.904)(5,6.723)(6,7.379)(7,7.903)(8,8.322)(9,8.658)(10,8.926)};
  \node[font=\sffamily\footnotesize, anchor=south west, text=bkBlue!80!black] at (axis cs:2.8,17) {左の水位 $h_1$};
  \node[font=\sffamily\footnotesize, anchor=north west, text=bkOrange!80!black] at (axis cs:3.2,3.0) {右の水位 $h_2$};
  \node[font=\sffamily\footnotesize, anchor=south west, text=bkGray] at (axis cs:5.6,10.2) {10\,cm};
\end{axis}
\end{tikzpicture}
\bkfigcaption{点は表計算（$\Delta t=1$ 分）、線は⑤で求める式の答え。2つの水位は、上と下から10\,cm に近づく。}
\end{bkfigure}

表計算の点は、線（ほんとうの答え）より少しだけ速く10\,cm に近づいています。表では差が毎分ちょうど0.8倍（$1-0.2=0.8$）になりますが、やってみようでは約0.82倍でした。6-2 と同じで、1分のあいだにも差は小さくなっていくのに、はじめの差のまま1分進めたからです。$\Delta t$ を小さくすると、ずれは小さくなります（10分後の $h_1$ は、$\Delta t=1$ で11.07、$\Delta t=0.5$ で11.22、$\Delta t=0.1$ で11.33。ほんとうは11.35）。

\bksubsection{4}{2回微分の式は、1回微分の式2つになる}

\label{bk:6-14:kaisuu}6-2 では、ばねの動きを、位置 $x$ と速度 $v$ の2つの列で計算しました。じつはこれも、連立微分方程式を解いていたのです。

ばね $m\ddot x=-kx$ は、2回微分が入った1つの式です。速度は位置の微分（$v=\dot x$）、加速度は速度の微分（$\ddot x=\dot v$）なので（5-5）、この式は次の2つの式の組と同じです。
\kvioframe[100mm]{$\dfrac{\dd x}{\dd t}=v$\qquad $\dfrac{\dd v}{\dd t}=-\dfrac{k}{m}\,x$}
1つめは「位置の変わる速さは速度」、2つめは「速度の変わる速さ（加速度）は、位置の $-k/m$ 倍」です。1つめの式に $v$、2つめの式に $x$ が入っているので、これは連立微分方程式です。
\begin{taisetsu}
2回微分の式1つは、1回微分の式2つ（位置と速度）の組に書きかえられる。
\end{taisetsu}

ある時刻の「位置と速度の組」がわかれば、その先の動きは1つに決まります（2回微分の式には、はじめの位置と速度の2つが必要でした。6-1）。この「組」を、その時刻の\textbf{状態}といいます。状態を、横に位置、たてに速度をとった平面の点で表すと、時間がたつにつれて点が動いていきます。

\begin{bkfigure}
\begin{tikzpicture}
\begin{axis}[bkplot, title={ダンパーなし}, title style={font=\sffamily\bfseries\footnotesize, yshift=5.5mm}, axis lines=box, grid=none, width=56mm, height=56mm, xmin=-0.13, xmax=0.13, ymin=-0.5, ymax=0.5, xtick={-0.1,0,0.1}, ytick={-0.4,0,0.4}, xlabel={位置 $x$（m）}, ylabel={速度 $v$（m/s）}, every axis x label/.style={at={(axis description cs:0.5,0)}, anchor=north, yshift=-5mm, font=\sffamily\footnotesize}, every axis y label/.style={at={(axis description cs:0,1)}, anchor=south west, yshift=1mm, font=\sffamily\footnotesize}, xticklabel style={/pgf/number format/fixed}, yticklabel style={/pgf/number format/fixed}]
  \addplot[bkthin, line width=0.5pt] coordinates {(-0.13,0)(0.13,0)};
  \addplot[bkthin, line width=0.5pt] coordinates {(0,-0.5)(0,0.5)};
  \addplot[bkcurve, domain=0:1.571, samples=120, variable=\t] ({0.1*cos(deg(4*\t))}, {-0.4*sin(deg(4*\t))});
  \addplot[only marks, mark=*, mark size=2pt, bkOrange] coordinates {(0.1,0)};
  \draw[bkBlue, line width=1pt, -{Stealth[length=5pt]}] (axis cs:0.0707,-0.283) -- (axis cs:0.06,-0.32);
\end{axis}
\end{tikzpicture}\hspace{10mm}%
\begin{tikzpicture}
\begin{axis}[bkplot, title={ダンパーあり}, title style={font=\sffamily\bfseries\footnotesize, yshift=5.5mm}, axis lines=box, grid=none, width=56mm, height=56mm, xmin=-0.13, xmax=0.13, ymin=-0.5, ymax=0.5, xtick={-0.1,0,0.1}, ytick={-0.4,0,0.4}, xlabel={位置 $x$（m）}, ylabel={速度 $v$（m/s）}, every axis x label/.style={at={(axis description cs:0.5,0)}, anchor=north, yshift=-5mm, font=\sffamily\footnotesize}, every axis y label/.style={at={(axis description cs:0,1)}, anchor=south west, yshift=1mm, font=\sffamily\footnotesize}, xticklabel style={/pgf/number format/fixed}, yticklabel style={/pgf/number format/fixed}]
  \addplot[bkthin, line width=0.5pt] coordinates {(-0.13,0)(0.13,0)};
  \addplot[bkthin, line width=0.5pt] coordinates {(0,-0.5)(0,0.5)};
  \addplot[bkcurve2, domain=0:9, samples=400, variable=\t] ({0.1*exp(-0.4*\t)*(cos(deg(3.98*\t))+0.1005*sin(deg(3.98*\t)))}, {-0.402*exp(-0.4*\t)*sin(deg(3.98*\t))});
  \addplot[only marks, mark=*, mark size=2pt, bkOrange] coordinates {(0.1,0)};
\end{axis}
\end{tikzpicture}
\bkfigcaption{状態の点 $(x,\ v)$ の動き（6-6・6-7 のばね：0.5\,kg、8\,N/m、はじめ $x=\qty{0.1}{m}$、$v=0$。オレンジの点がスタート）。ダンパーなし：同じだ円を回り続ける。ダンパーあり（$c=0.4$\,N·s/m）：うずをまいて原点（止まった状態）へ近づく。}\label{bk:6-14:fig-joutai}
\end{bkfigure}

図\ref{bk:6-14:fig-joutai}の「ダンパーなし」では、点は時計回りに同じだ円を回り続けます。ばねのゆれが、いつまでも同じ大きさでくり返すからです。「ダンパーあり」では、ダンパーがエネルギーをうばうので、点はうずをまいて原点（位置0・速度0、つまり止まった状態）に近づきます。この「状態の点の動き」で見る考え方は、第9巻（9-10）の状態方程式で、機械やロボットを制御するときの基本になります。

\kviohead{2つの台車を、4つの列で進める}
2つの台車は、位置 $x_1$、$x_2$ と速度 $v_1$、$v_2$ の4つで状態が決まります。6-2 の型（速度を先に進める）で、4つの列を同時に進めます。加速度は、②の運動方程式を $m$ でわったものです。$m=\qty{1}{kg}$、$k=\qty{100}{N/m}$、$K=\qty{10}{N/m}$ なら、
\kvioline{$a_1=\dfrac{-kx_1+K(x_2-x_1)}{m}=-110x_1+10x_2$}
\kvioline{$a_2=\dfrac{-kx_2-K(x_2-x_1)}{m}=10x_1-110x_2$}
（$-kx_1-Kx_1=-(100+10)x_1=-110x_1$ のように、同じ文字の項をまとめました。）

\begin{bkcenter}
\bksheet{7}{15}{%
1 & $t$ & $x_1$ & $x_2$ & $v_1$ & $v_2$ & $a_1$ & $a_2$ \\
2 & 0 & 0.1 & 0 & 0 & 0 & $-11$ & 1 \\
3 & 0.01 & 0.0989 & 0.0001 & $-0.11$ & 0.01 & $-10.878$ & 0.978 \\
4 & 0.02 & 0.096712 & 0.000298 & $-0.21878$ & 0.01978 & $-10.635$ & 0.934 \\}
\bktabcaption{$\Delta t=\qty{0.01}{s}$ で進めた表のはじめ（台車1を\qty{0.1}{m}引いてはなした）。}
\end{bkcenter}
表の数は、次の順で計算します。
\begin{enumerate}[label={\bkcircled[\bkTT]{\arabic*}}, leftmargin=2.4\zw, itemsep=1pt, topsep=2pt]
\item 2行めの加速度：$a_1=-110\times0.1+10\times0=-11$\newline $a_2=10\times0.1-110\times0=1$
\item 3行めの速度：$v_1=0+(-11)\times0.01=-0.11$\newline $v_2=0+1\times0.01=0.01$
\item 3行めの位置（新しい速度で）：$x_1=0.1+(-0.11)\times0.01=0.0989$\newline $x_2=0+0.01\times0.01=0.0001$
\end{enumerate}
台車2は、はじめは止まっていますが、つなぎのばねに引かれて、少しずつ動き出します。図\ref{bk:6-14:fig-uneri}は、この表を $\Delta t=\qty{0.001}{s}$ で7000行計算してかいたものです。

\bkkokomade[表計算で、もう未来は計算できます。この先は式で解きます。]

\bkstage{shiki}{式でわかる}{組み合わせて、ほどく}{\bklevel{5}}

\bksubsection{5}{和と差に分ける}

\label{bk:6-14:wasa}やってみようで気づいた「合計は一定」「差は同じ割合でへる」を、式で確かめます。2つの式をもう一度書きます。
\kvioline{$\dfrac{\dd h_1}{\dd t}=-0.1(h_1-h_2)$\qquad $\dfrac{\dd h_2}{\dd t}=0.1(h_1-h_2)$}

\kviohead{2つの式をたす：合計 $S=h_1+h_2$}
たし算の微分は、微分のたし算です（5-4：$(f+g)'=f'+g'$）。2つの式の左辺どうし、右辺どうしをたします。右辺は $-0.1(h_1-h_2)+0.1(h_1-h_2)$ で、打ち消しあって0になります。
\begin{bkhenkei}
\frac{\dd h_1}{\dd t}+\frac{\dd h_2}{\dd t}=0 & 2つの式をたした \\
\frac{\dd (h_1+h_2)}{\dd t}=0 & 左辺は、たし算の微分（5-4） \\
\frac{\dd S}{\dd t}=0 & $S=h_1+h_2$ と置いた \\
\end{bkhenkei}
変わる速さが0なので、$S$ は一定です。はじめの値 $S=20+0=20$ のままです。左で下がった分だけ、右で上がるからです。

\kviohead{2つの式をひく：差 $D=h_1-h_2$}
こんどは、左の式から右の式をひきます。右辺は、次のようになります。
\kvioline{$-0.1(h_1-h_2)-0.1(h_1-h_2)=-0.2(h_1-h_2)$}
\begin{bkhenkei}
\frac{\dd h_1}{\dd t}-\frac{\dd h_2}{\dd t}=-0.2(h_1-h_2) & 左の式から右の式をひいた \\
\frac{\dd (h_1-h_2)}{\dd t}=-0.2(h_1-h_2) & 左辺は、ひき算の微分（5-4） \\
\frac{\dd D}{\dd t}=-0.2D & $D=h_1-h_2$ と置いた \\
D=20e^{-0.2t} & 6-4 の形。はじめの差は20 \\
\end{bkhenkei}
$D$ だけの式になりました。「自分に比例してへる」形（6-4）なので、答えは $D=D_0e^{-0.2t}$、はじめの差 $D_0=20-0=20$ です。時定数は $\tau=1\div0.2=5$ 分。やってみようの「毎分約0.82倍」は、$e^{-0.2}=0.8187$ だったのです。

\kviohead{もとの $h_1$、$h_2$ にもどす}
$S=h_1+h_2$ と $D=h_1-h_2$ から、$h_1$ と $h_2$ を求めます。3-6 の加減法と同じです。
\begin{bkhenkei}
S+D=2h_1 & たすと $h_2$ が消える \\
h_1=\frac{S+D}{2}=10+10e^{-0.2t} & 2でわって、$S$ と $D$ を入れた \\
S-D=2h_2 & ひくと $h_1$ が消える \\
h_2=\frac{S-D}{2}=10-10e^{-0.2t} & 2でわって、$S$ と $D$ を入れた \\
\end{bkhenkei}

\kviohead{代入して確かめる}
$h_1=10+10e^{-0.2t}$ を微分すると（5-14）、
\kvioline{$\dfrac{\dd h_1}{\dd t}=10\times(-0.2)e^{-0.2t}=-2e^{-0.2t}$}
右辺は $-0.1(h_1-h_2)=-0.1\times20e^{-0.2t}=-2e^{-0.2t}$。同じなので、たしかに答えです。$h_2$ も同じように確かめられます。

\kviohead{\bkhatenaref の答え}
左の水位は $h_1=10+10e^{-0.2t}$、右の水位は $h_2=10-10e^{-0.2t}$（cm、分）。どちらも最後は\qty{10}{cm}になります。2つの水位の差は、5分（時定数）ごとに $e^{-1}\approx0.37$ 倍、$5\ln2\approx3.5$ 分ごとに半分になります。

\begin{bkfigure}
\begin{tikzpicture}
\begin{axis}[bkplot, width=96mm, height=46mm, xmin=0, xmax=15.8, ymin=0, ymax=23, xtick={5,10,15}, ytick={10,20}, xlabel={$t$（分）}, ylabel={cm}]
  \addplot[bkcurve3, domain=0:15.3] {20};
  \addplot[bkPurple, line width=1.7pt, domain=0:15.3, samples=60] {20*exp(-0.2*x)};
  \node[font=\sffamily\footnotesize, anchor=south, text=bkGreen!70!black] at (axis cs:11,20.3) {合計 $S$（一定）};
  \node[font=\sffamily\footnotesize, anchor=south west, text=bkPurple] at (axis cs:6.3,9) {差 $D=20e^{-0.2t}$};
\end{axis}
\end{tikzpicture}
\bkfigcaption{合計と差に組みかえると、片方は変わらず、もう片方は6-4の形でへるだけ。}
\end{bkfigure}

\begin{seikaku}
和と差に分けると2つの式がばらばらになったのは、2つの水そうが同じ形で、式が左右で同じ形（左と右を入れかえても同じ式）だったからです。形がちがう水そうでは、ただの和ではうまくいきません。次の小見出しのように、底面積をかけた「重みつきの和」を使います。
\end{seikaku}

\kviohead{大きさのちがう水そうでは}
左の水そうの底面積を\qty{100}{cm^2}、右を\qty{300}{cm^2}にします。管は同じで、1分に $10(h_1-h_2)$\,cm³ 流れます。はじめは左が\qty{20}{cm}、右が空です。
\kvioline{$\dfrac{\dd h_1}{\dd t}=-\dfrac{10(h_1-h_2)}{100}=-0.1(h_1-h_2)$}
\kvioline{$\dfrac{\dd h_2}{\dd t}=\dfrac{10(h_1-h_2)}{300}=\dfrac{1}{30}(h_1-h_2)$}
右は底面積が3倍なので、同じ水が入っても、水位は3分の1しか上がりません。だから $h_1+h_2$ は一定ではありません。一定なのは、水の体積 $100h_1+300h_2$ です。確かめます。
\kvioline{$\dfrac{\dd}{\dd t}(100h_1+300h_2)=100\times\bigl(-0.1(h_1-h_2)\bigr)+300\times\dfrac{1}{30}(h_1-h_2)$}
\kvioline{$=-10(h_1-h_2)+10(h_1-h_2)=0$}
このように、底面積の重みをかけてたしたものを、\textbf{重みつきの和}といいます。はじめの体積は $100\times20+300\times0=2000$\,cm³ です。差 $D=h_1-h_2$ のほうは、
\kvioline{$\dfrac{\dd D}{\dd t}=-0.1D-\dfrac{1}{30}D=-\dfrac{3+1}{30}D=-\dfrac{2}{15}D$}
で、やはり6-4の形です。時定数は $\tau=15\div2=7.5$ 分、$D=20e^{-t/7.5}$ です。最後は2つの水位がそろうので、$100h+300h=2000$ より $h=\qty{5}{cm}$。体積の式と差の式から $h_1$、$h_2$ を出します（3-6 の代入法：$h_1=h_2+D$ を体積の式に入れる）。
\kvioline{$100(h_2+D)+300h_2=2000$\quad より\quad $h_2=5-\dfrac{D}{4}=5-5e^{-t/7.5}$}
\kvioline{$h_1=h_2+D=5+15e^{-t/7.5}$}
左は\qty{15}{cm}下がり、右は\qty{5}{cm}しか上がりません。右の底面積が3倍だからです。

\begin{bkfigure}
\begin{tikzpicture}
\begin{axis}[bkplot, width=100mm, height=54mm, xmin=0, xmax=31.5, ymin=0, ymax=22, xtick={10,20,30}, ytick={5,10,15,20}, xlabel={$t$（分）}, ylabel={水位（cm）}]
  \addplot[bkthin, dashed, domain=0:30.5] {5};
  \addplot[bkcurve, domain=0:30.5, samples=60] {5+15*exp(-x/7.5)};
  \addplot[bkcurve2, domain=0:30.5, samples=60] {5-5*exp(-x/7.5)};
  \node[font=\sffamily\footnotesize, anchor=south west, text=bkBlue!80!black] at (axis cs:6,12.5) {$h_1=5+15e^{-t/7.5}$};
  \node[font=\sffamily\footnotesize, anchor=north west, text=bkOrange!80!black] at (axis cs:12,3.7) {$h_2=5-5e^{-t/7.5}$};
\end{axis}
\end{tikzpicture}
\bkfigcaption{底面積が100\,cm² と300\,cm² の水そう。2つの水位は5\,cm にそろう。}
\end{bkfigure}

\bksubsection{6}{2つのおもり：振動モード}

\label{bk:6-14:mode}②の台車の式にも、和と差を使います。
\kvioline{$m\ddot x_1=-kx_1+K(x_2-x_1)$\qquad $m\ddot x_2=-kx_2-K(x_2-x_1)$}

\kviohead{和 $u=x_1+x_2$}
2つの式をたします。
\kvioline{$m(\ddot x_1+\ddot x_2)=-kx_1+K(x_2-x_1)-kx_2-K(x_2-x_1)$}
\begin{bkhenkei}
m(\ddot x_1+\ddot x_2)=-k(x_1+x_2) & つなぎのばねの力は打ち消しあう \\
m\ddot u=-ku & $u=x_1+x_2$ と置いた \\
\ddot u=-\frac{k}{m}u & 両辺を $m$ でわった。6-6 の形 \\
\end{bkhenkei}
\kviohead{差 $w=x_1-x_2$}
左の式から右の式をひきます。
\kvioline{$m(\ddot x_1-\ddot x_2)=-kx_1+K(x_2-x_1)+kx_2+K(x_2-x_1)$}
\begin{bkhenkei}
m(\ddot x_1-\ddot x_2)=-(k+2K)(x_1-x_2) & $K(x_2-x_1)=-K(x_1-x_2)$ が2つ分 \\
m\ddot w=-(k+2K)w & $w=x_1-x_2$ と置いた \\
\ddot w=-\frac{k+2K}{m}w & 両辺を $m$ でわった \\
\end{bkhenkei}
どちらも、6-6 の「2回微分すると自分の $-\omega^2$ 倍」の形です。だから、
\kvioframe[112mm]{和 $u$ は\quad $\omega_1=\sqrt{\dfrac{k}{m}}$\quad でゆれる\\[2mm] 差 $w$ は\quad $\omega_2=\sqrt{\dfrac{k+2K}{m}}$\quad でゆれる}

\kviohead{2つの特別な動き方}
式の意味を、図で考えます。
\begin{bkfigure}
\begin{tikzpicture}[x=1mm, y=1mm]
  \begin{scope}
  \draw[bkInk, line width=0.8pt] (0,0) -- (124,0);
  \kviowall{(2,0)}{14}{1}\kviowall{(122,0)}{14}{-1}
  \kviospring{(2,7)}{(36,7)}\kviocart{(36,0)}{1}{bkBlue}\kviospring{(48,7)}{(86,7)}\kviocart{(86,0)}{2}{bkOrange}\kviospring{(98,7)}{(122,7)}
  \draw[bkRed, line width=1.3pt, -{Stealth[length=4pt]}] (38,13) -- (48,13);
  \draw[bkRed, line width=1.3pt, -{Stealth[length=4pt]}] (88,13) -- (98,13);
  \node[font=\sffamily\footnotesize, anchor=south] at (67,10.5) {のびない};
  \node[font=\sffamily\bfseries\small, anchor=west, text=\bkTT] at (0,21) {1次モード：いっしょに動く\quad $x_1=x_2$、$\omega_1=\sqrt{k/m}$};
  \end{scope}
  \begin{scope}[yshift=-34mm]
  \draw[bkInk, line width=0.8pt] (0,0) -- (124,0);
  \kviowall{(2,0)}{14}{1}\kviowall{(122,0)}{14}{-1}
  \kviospring{(2,7)}{(36,7)}\kviocart{(36,0)}{1}{bkBlue}\kviospring{(48,7)}{(86,7)}\kviocart{(86,0)}{2}{bkOrange}\kviospring{(98,7)}{(122,7)}
  \draw[bkRed, line width=1.3pt, -{Stealth[length=4pt]}] (46,13) -- (36,13);
  \draw[bkRed, line width=1.3pt, -{Stealth[length=4pt]}] (88,13) -- (98,13);
  \node[font=\sffamily\footnotesize, anchor=south] at (67,10.5) {2倍のびる};
  \node[font=\sffamily\bfseries\small, anchor=west, text=\bkTT] at (0,21) {2次モード：反対に動く\quad $x_1=-x_2$、$\omega_2=\sqrt{(k+2K)/m}$};
  \end{scope}
\end{tikzpicture}
\bkfigcaption{2つの台車がいっしょに動くと、つなぎのばねはのびもちぢみもしない。反対に動くと、つなぎのばねは片方の台車のずれの2倍のびちぢみする。}
\end{bkfigure}

\begin{itemize}[leftmargin=1.2\zw, itemsep=3pt, topsep=2pt, label={\color{\bkTT}●}]
\item \textbf{いっしょに動く（$x_1=x_2$）}\par\noindent 2つの台車が同じだけ右へ動くと、つなぎのばねの長さは変わりません。だから、それぞれの台車は、外側のばね $k$ だけに引かれます。台車1つとばね1本のときと同じなので、$\omega_1=\sqrt{k/m}$ でゆれます。
\item \textbf{反対に動く（$x_1=-x_2$）}\par\noindent 台車1が右へ $x_1$、台車2が左へ $x_1$ 動くと、つなぎのばねは $2x_1$ ちぢみます。台車1は、外側のばねから $kx_1$、つなぎのばねから $2Kx_1$ の力で押しもどされます。ばねが $k+2K$ に強くなったのと同じなので、速くゆれます。
\end{itemize}
このように、全部の部分が同じ角振動数で、そろってゆれる動き方を\textbf{振動モード}（しんどうモード）といいます。\bkidx{しんどうもーど}{振動モード}角振動数の小さいほうから、1次モード、2次モードとよびます。

2つの台車を同じ向きに同じだけずらしてはなすと、差 $w$ ははじめから0で、ずっと0のままです。だから、いつまでも1次モードだけでゆれます。反対向きに同じだけずらすと、和 $u$ が0のままなので、2次モードだけでゆれます。

\kviohead{どんな動きも、2つのモードのたし算}
台車1だけを $X$ 引いて、2つとも静かにはなしたときを考えます。はじめは $x_1=X$、$x_2=0$ なので、$u=X+0=X$、$w=X-0=X$。速度は0です。6-6 と同じように、はじめの位置が $X$ で速度が0なら、答えは $\cos$ の形です。
\kvioline{$u=X\cos\omega_1t$\qquad $w=X\cos\omega_2t$}
$x_1=(u+w)\div2$、$x_2=(u-w)\div2$ にもどすと（⑤と同じ加減法）、
\kvioframe[100mm]{$x_1=\dfrac X2(\cos\omega_1t+\cos\omega_2t)$\qquad $x_2=\dfrac X2(\cos\omega_1t-\cos\omega_2t)$}
台車1の動きも、台車2の動きも、2つのモードのゆれのたし算です。

②の数（$m=1$、$k=100$、$K=10$、$X=0.1$）では、$\omega_1=\sqrt{100}=10$、$\omega_2=\sqrt{120}\approx10.95$（rad/s）です。2つの角振動数が近いので、4-9 のうなりになります。加法定理（4-9）の $\cos(P\pm Q)=\cos P\cos Q\mp\sin P\sin Q$ を、たしたり、ひいたりすると、
\kvioline{$\cos(P+Q)+\cos(P-Q)=2\cos P\cos Q$}
\kvioline{$\cos(P+Q)-\cos(P-Q)=-2\sin P\sin Q$}
$P+Q=\omega_2t$、$P-Q=\omega_1t$ になるように、$P=\dfrac{\omega_1+\omega_2}{2}t$、$Q=\dfrac{\omega_2-\omega_1}{2}t$ とします。すると、
\kvioline{$x_1=X\cos\dfrac{\omega_2-\omega_1}{2}t\cdot\cos\dfrac{\omega_1+\omega_2}{2}t$}
\kvioline{$x_2=-X\sin\dfrac{\omega_2-\omega_1}{2}t\cdot\sin\dfrac{\omega_1+\omega_2}{2}t$}
$(\omega_1+\omega_2)\div2\approx10.48$ の速いゆれの大きさが、$(\omega_2-\omega_1)\div2\approx0.477$ のゆっくりした $\cos$ と $\sin$ で変わります。$\cos0.477t$ が0になるのは $0.477t=\pi\div2$、つまり $t\approx\qty{3.3}{s}$ のときです。このとき台車1は止まり、台車2のゆれの大きさは $X$ です。図\ref{bk:6-14:fig-uneri}のふしぎな動きは、こうして式で説明できました。

3本とも同じばね（$K=k=\qty{100}{N/m}$）なら、$\omega_1=10$、$\omega_2=\sqrt{300}\approx17.3$（rad/s）です。周期は $2\pi\div10\approx\qty{0.63}{s}$ と $2\pi\div17.3\approx\qty{0.36}{s}$。2つの角振動数がはなれているので、うなりではなく、こみいった形のゆれになります。それでも、2つのモードのたし算であることは同じです。

\kviohead{係数を表に並べると}
②の式を $m$ でわって、$x_1$ と $x_2$ の係数を表に並べてみます。
\begin{bkcenter}
\begin{bktable}{colspec={Q[c,m,wd=20mm] Q[c,m,wd=32mm] Q[c,m,wd=32mm]}, rowsep=4pt, row{1}={font=\sffamily\small}}
 & $x_1$ の係数 & $x_2$ の係数 \\
$\ddot x_1=$ & $-\dfrac{k+K}{m}$ & $\dfrac{K}{m}$ \\
$\ddot x_2=$ & $\dfrac{K}{m}$ & $-\dfrac{k+K}{m}$ \\
\end{bktable}
\bktabcaption{2つの式の係数を、表の形に並べた。}
\end{bkcenter}
$x_1$ と $x_2$ を「1対1でたす（和）」「1対1でひく（差）」という組み合わせにしたら、2つの式がばらばらにほどけました。こういう「ほどける組み合わせ」を、係数の表（行列）の計算で見つける方法が、第7巻の\textbf{固有ベクトル}（7-8）と\textbf{対角化}（7-9）です。おもりが3つ、4つ…とふえても、表の計算でモードを見つけられるようになります。

\bksubsection{7}{消去法：1つの式にまとめる}

\label{bk:6-14:shokyo}和と差のような、うまい組み合わせがいつも見つかるとはかぎりません。そのときは、3-6 の代入法と同じように、文字を1つ消します。1つの式を「$h_2=\cdots$」の形にして、もう1つの式に入れるのです。この方法を\textbf{消去法}（しょうきょほう）といいます。\bkidx{しょうきょほう}{消去法}

\kviohead{水そうの式で}
\begin{bkhenkei}
\frac{\dd h_1}{\dd t}=-0.1(h_1-h_2) & 左の水そうの式 \\
10\dot h_1=-h_1+h_2 & 両辺に10をかけた \\
h_2=h_1+10\dot h_1 & $h_2=\cdots$ の形にした \\
\dot h_1+10\ddot h_1=0.1\bigl(h_1-(h_1+10\dot h_1)\bigr) & 右の水そうの式に入れた \\
\dot h_1+10\ddot h_1=-\dot h_1 & 右辺をかっこの中から計算した \\
10\ddot h_1+2\dot h_1=0 & $h_1$ だけの、2回微分の式 \\
\end{bkhenkei}
4行めは、右の水そうの式 $\dot h_2=0.1(h_1-h_2)$ の $h_2$ に、3行めの $h_1+10\dot h_1$ を入れたものです。左辺は、$h_2=h_1+10\dot h_1$ を微分した $\dot h_2=\dot h_1+10\ddot h_1$ です。2つの1回微分の式から、$h_1$ だけの2回微分の式が1つできました。

これは 6-12 の形です。$h_1=e^{st}$ と置くと、特性方程式は
\kvioline{$10s^2+2s=0$\qquad $s(10s+2)=0$\qquad $s=0,\ -0.2$}
$e^{0t}=1$ なので、答えは $h_1=A+Be^{-0.2t}$ です。はじめの値 $h_1(0)=20$ と、はじめの変わる速さ $\dot h_1(0)=-0.1\times(20-0)=-2$ で $A$、$B$ を決めます。
\kvioline{$A+B=20$、\quad $-0.2B=-2$\quad より\quad $B=10$、\ $A=10$}
⑤と同じ $h_1=10+10e^{-0.2t}$ になりました。$s=0$ の解は「変わらない部分」で、合計が一定であることを表しています。

\kviohead{直列の水そう：上から下へ流れ落ちる}
こんどは、同じ形の水そうを、階段のように上下に置きます。上の水そうの水は、管を通って下の水そうに落ち、下の水そうの水は、同じような管から外に出ていきます。どちらの管も、流れる量はその水そうの水位に比例するとします。
\begin{bkfigure}
\begin{minipage}[c]{0.36\linewidth}
\centering
\begin{tikzpicture}[x=0.85mm, y=0.85mm]
  \kviotank{(0,22)}{18}{20}{16}
  \kviotank{(14,0)}{18}{20}{6}
  \draw[bkInk!75, line width=0.7pt] (18,23.2) -- (23,23.2) -- (23,19);
  \draw[bkInk!75, line width=0.7pt] (32,1.2) -- (40,1.2);
  \draw[bkBlue!80!black, line width=1pt, -{Stealth[length=3pt]}] (23,18.5) -- (23,13);
  \draw[bkBlue!80!black, line width=1pt, -{Stealth[length=3pt]}] (40.5,1.2) -- (45,1.2);
  \node[font=\sffamily\footnotesize, anchor=east] at (-1,30) {$h_1$};
  \node[font=\sffamily\footnotesize, anchor=west] at (33,8) {$h_2$};
\end{tikzpicture}
\end{minipage}\hfill
\begin{minipage}[c]{0.6\linewidth}
\begin{tikzpicture}
\begin{axis}[bkplot, width=68mm, height=46mm, xmin=0, xmax=21, ymin=0, ymax=44, xtick={4,8,12,16,20}, ytick={10,20,30,40}, xlabel={$t$（分）}, ylabel={水位（cm）}]
  \addplot[bkcurve, domain=0:20.5, samples=60] {40*exp(-x/4)};
  \addplot[bkcurve2, domain=0:20.5, samples=80] {10*x*exp(-x/4)};
  \addplot[only marks, mark=*, mark size=1.8pt, bkOrange!80!black] coordinates {(4,14.715)};
  \node[font=\sffamily\footnotesize, anchor=south west, text=bkBlue!80!black] at (axis cs:1.6,28) {上 $h_1$};
  \node[font=\sffamily\footnotesize, anchor=south west, text=bkOrange!80!black] at (axis cs:8.5,14.5) {下 $h_2$};
\end{axis}
\end{tikzpicture}
\end{minipage}
\bkfigcaption{直列の水そう（時定数4分、はじめ上が40\,cm、下が空）。下の水位は、ふえてから、へる。いちばん高くなるのは4分後（点）。}
\end{bkfigure}

それぞれの水そうだけで流れ出すときの時定数を $\tau$ とすると（6-4）、
\kvioframe[112mm]{$\dfrac{\dd h_1}{\dd t}=-\dfrac{h_1}{\tau}$\qquad $\dfrac{\dd h_2}{\dd t}=\dfrac{h_1}{\tau}-\dfrac{h_2}{\tau}$}
下の水そうには、上から $h_1/\tau$ の割合で水位がふえるように水が入り、自分の管から $h_2/\tau$ の割合で出ていきます。上の式には $h_1$ しかないので、6-4 で $h_1=h_0e^{-t/\tau}$（$h_0$ ははじめの深さ）とすぐに解けます。下の式は、入る量が時間とともに変わるので、6-4 の形ではありません。そこで消去法を使います。
\begin{bkhenkei}
h_1=\tau\dot h_2+h_2 & 下の式の両辺に $\tau$ をかけて、$h_1=\cdots$ にした \\
\tau\ddot h_2+\dot h_2=-\frac{\tau\dot h_2+h_2}{\tau} & 上の式 $\dot h_1=-h_1/\tau$ に入れた \\
\tau^2\ddot h_2+2\tau\dot h_2+h_2=0 & 両辺に $\tau$ をかけて、左辺に集めた \\
\tau^2s^2+2\tau s+1=0 & $h_2=e^{st}$ と置いた（6-12） \\
(\tau s+1)^2=0 & 2乗の展開の形（3-8） \\
s=-\frac{1}{\tau} & 重解 \\
\end{bkhenkei}
重解なので、6-12 のとおり、答えは $h_2=(A+Bt)e^{-t/\tau}$ です。$A$ と $B$ は、はじめの状態で決めます。
\begin{enumerate}[label={\bkcircled[\bkTT]{\arabic*}}, leftmargin=2.4\zw, itemsep=1pt, topsep=2pt]
\item はじめは下が空なので、$h_2(0)=A=0$。
\item はじめの変わる速さは、下の式に $h_1=h_0$、$h_2=0$ を入れて $\dot h_2(0)=h_0/\tau$。
\item $h_2=Bte^{-t/\tau}$ を積の微分（5-6）で微分すると $\dot h_2=B\Bigl(1-\dfrac{t}{\tau}\Bigr)e^{-t/\tau}$。$t=0$ では $B$ なので、$B=h_0/\tau$。
\end{enumerate}
\kvioframe[80mm]{$h_2=h_0\,\dfrac{t}{\tau}\,e^{-t/\tau}$}
6-12 では、$te^{st}$ は「重解のときに足りない2つめの答え」として出てきました。ここでは、それが主役です。下の水位がいちばん高くなるのは、微分が0になる時刻です（5-7）。③の $\dot h_2$ が0になるのは $1-t/\tau=0$、つまり $t=\tau$ のときで、そのときの水位は $h_0e^{-1}\approx0.37h_0$ です。

\kviohead{時定数がちがうとき：親と子}
2つの水そうの時定数がちがうとき（上が $\tau_1$、下が $\tau_2$）は、量を $N_1$、$N_2$ と書くと、
\kvioline{$\dfrac{\dd N_1}{\dd t}=-\dfrac{N_1}{\tau_1}$\qquad $\dfrac{\dd N_2}{\dd t}=\dfrac{N_1}{\tau_1}-\dfrac{N_2}{\tau_2}$}
となります。これは、放射性物質の「親」がこわれて「子」になり、子もまたこわれていくときの式と同じです。原子の数がへる速さは、そのときの数に比例します（6-4）。親の数 $N_1$ は親のこわれ方だけで決まり、子の数 $N_2$ は、親から生まれる分だけふえ、自分がこわれる分だけへります。この式は、1910年にベイトマンが一般の場合まで解きました。

同じように消去法を使うと、特性方程式は $(\tau_1s+1)(\tau_2s+1)=0$ になり（例題8でくわしく計算します）、解は $s=-1/\tau_1$ と $s=-1/\tau_2$ の2つです。だから子の数は、
\kvioline{$N_2=Ae^{-t/\tau_1}+Be^{-t/\tau_2}$}
の形になります。

\begin{bkfigure}
\begin{tikzpicture}
\begin{axis}[bkplot, width=100mm, height=48mm, xmin=0, xmax=8.4, ymin=0, ymax=1100, xtick={2,4,6,8}, ytick={250,500,750,1000}, xlabel={$t$（h）}, ylabel={数}]
  \addplot[bkcurve, domain=0:8.1, samples=60] {1000*exp(-x/2)};
  \addplot[bkcurve2, domain=0:8.1, samples=80] {1000*(exp(-x/2)-exp(-x))};
  \addplot[only marks, mark=*, mark size=1.8pt, bkOrange!80!black] coordinates {(1.386,250)};
  \node[font=\sffamily\footnotesize, anchor=south west, text=bkBlue!80!black] at (axis cs:2.2,420) {親 $N_1=1000e^{-t/2}$};
  \node[font=\sffamily\footnotesize, anchor=south west, text=bkOrange!80!black] at (axis cs:3.6,190) {子 $N_2=1000(e^{-t/2}-e^{-t})$};
\end{axis}
\end{tikzpicture}
\bkfigcaption{例題8の親と子（$\tau_1=\qty{2}{h}$、$\tau_2=\qty{1}{h}$、はじめ親1000、子0）。子は、ふえてからへる。直列の水そうと同じ形。}
\end{bkfigure}
2つの時定数が近づいていくと、この答えは直列の水そうの $te^{-t/\tau}$ の形に近づいていきます。

\bksubsection{8}{熱と電気でも同じ形}

\label{bk:6-14:netsu}\kviohead{2つの物の温度が近づく}
熱いお湯を入れた金属のコップを、冷たい水を入れた入れ物にひたします。外へ熱は逃げないとします（入れ物を発ぽうスチロールの箱に入れる）。熱は、温度の高いほうから低いほうへ流れます。1秒に流れる熱（熱流、W）は、温度差に比例し、熱抵抗 $R_{\mathrm{th}}$〔K/W〕に反比例します（6-4）。
\kvioline{熱流 $P=\dfrac{T_1-T_2}{R_{\mathrm{th}}}$\quad 〔W〕}
お湯（熱容量 $C_1$〔J/K〕）は1秒に $P$ の熱を失い、水（熱容量 $C_2$）は同じ $P$ を受けとります。熱容量は「1\,K（1\,℃）上げるのに必要な熱」なので、温度の変わる速さは「熱の出入り÷熱容量」です。
\kvioframe[112mm]{\small $C_1\dfrac{\dd T_1}{\dd t}=-\dfrac{T_1-T_2}{R_{\mathrm{th}}}$\qquad $C_2\dfrac{\dd T_2}{\dd t}=\dfrac{T_1-T_2}{R_{\mathrm{th}}}$}
大きさのちがう水そうと、まったく同じ形です。2つの式をたすと、
\kvioline{$\dfrac{\dd}{\dd t}(C_1T_1+C_2T_2)=0$}
熱容量をかけた重みつきの和（熱の量）が一定です。最後は2つの温度がそろうので、最後の温度は
\kvioline{$T_\infty=\dfrac{C_1T_1(0)+C_2T_2(0)}{C_1+C_2}$}
です。これは、中学の理科で「お湯と水をまぜたときの温度」を求めた式と同じです（まぜても、金属のかべごしでも、熱の合計は変わらないから）。差 $D=T_1-T_2$ は、
\kvioline{$\dfrac{\dd D}{\dd t}=\dfrac{\dd T_1}{\dd t}-\dfrac{\dd T_2}{\dd t}=-\dfrac{D}{R_{\mathrm{th}}C_1}-\dfrac{D}{R_{\mathrm{th}}C_2}$}
\kvioline{$=-\dfrac{1}{R_{\mathrm{th}}}\Bigl(\dfrac{1}{C_1}+\dfrac{1}{C_2}\Bigr)D$}
で、6-4 の形でへります。2つの温度は、上と下から $T_\infty$ に近づきます（例題4）。

\kviohead{2つのコンデンサーをつなぐ}
\begin{bkfigure}
\begin{minipage}[c]{0.42\linewidth}
\centering
\begin{tikzpicture}[x=0.76mm, y=0.76mm]
  \draw[bkInk, line width=0.9pt] (0,0) -- (0,9.6) (0,12.4) -- (0,22) -- (14,22);
  \draw[bkInk, line width=1.4pt] (-4,9.6) -- (4,9.6) (-4,12.4) -- (4,12.4);
  \draw[bkInk, line width=0.9pt] (14,19.5) rectangle (28,24.5);
  \draw[bkInk, line width=0.9pt] (28,22) -- (32,22);
  \draw[bkInk, line width=0.9pt] (32,22) -- (37,25);
  \fill[bkInk] (32,22) circle (0.7) (38,22) circle (0.7);
  \draw[bkInk, line width=0.9pt] (38,22) -- (46,22) -- (46,12.4) (46,9.6) -- (46,0) -- (0,0);
  \draw[bkInk, line width=1.4pt] (42,9.6) -- (50,9.6) (42,12.4) -- (50,12.4);
  \node[font=\sffamily\footnotesize, anchor=east] at (-5,11) {$C_1$};
  \node[font=\sffamily\footnotesize, anchor=west] at (51,11) {$C_2$};
  \node[font=\sffamily\footnotesize, anchor=south] at (21,25.5) {$R$};
  \node[font=\sffamily\footnotesize, anchor=south] at (36,26) {スイッチ};
  \draw[bkOrange!85!black, line width=1pt, -{Stealth[length=3pt]}] (16,16) -- (26,16);
  \node[font=\sffamily\footnotesize, anchor=north, text=bkOrange!80!black] at (21,14.5) {電流 $i$};
\end{tikzpicture}
\end{minipage}\hfill
\begin{minipage}[c]{0.54\linewidth}
\begin{tikzpicture}
\begin{axis}[bkplot, width=56mm, height=42mm, xmin=0, xmax=5.3, ymin=0, ymax=11.5, xtick={1,2,3,4,5}, ytick={5,10}, xlabel={$t$（s）}, ylabel={電圧（V）}]
  \addplot[bkthin, dashed, domain=0:5.1] {5};
  \addplot[bkcurve, domain=0:5.1, samples=60] {5+5*exp(-x)};
  \addplot[bkcurve2, domain=0:5.1, samples=60] {5-5*exp(-x)};
  \node[font=\sffamily\footnotesize, anchor=south west, text=bkBlue!80!black] at (axis cs:0.9,7.4) {$V_1$};
  \node[font=\sffamily\footnotesize, anchor=north west, text=bkOrange!80!black] at (axis cs:0.9,2.6) {$V_2$};
\end{axis}
\end{tikzpicture}
\end{minipage}
\bkfigcaption{\qty{10}{V}に充電したコンデンサー（\qty{100}{\micro F}）と、空のコンデンサー（\qty{100}{\micro F}）を、\qty{20}{k\ohm}の抵抗でつなぐ。2つの電圧は\qty{5}{V}に近づく。}
\end{bkfigure}
コンデンサーは、電気をためる部品です（6-4）。ためた電気の量 $q$〔C〕と電圧 $V$〔V〕には $q=CV$ のきまりがあり、流れこむ電流 $i$〔A〕は $q$ の変わる速さです（$i=\dd q/\dd t$）。2つのコンデンサーを抵抗 $R$ でつなぐと、オームの法則で $i=(V_1-V_2)/R$ の電流が流れます。左のコンデンサーからは電気が出ていき、右には入るので、
\kvioline{$C_1\dfrac{\dd V_1}{\dd t}=-\dfrac{V_1-V_2}{R}$\qquad $C_2\dfrac{\dd V_2}{\dd t}=\dfrac{V_1-V_2}{R}$}
熱の式の $T$ を $V$ に、$R_{\mathrm{th}}$ を $R$ にかえただけです。電気の量 $C_1V_1+C_2V_2$ が一定で、電圧の差が6-4の形でへります。図のように $C_1=C_2=\qty{100}{\micro F}$、$R=\qty{20}{k\ohm}$ なら、差がへる速さは
\kvioline{$\dfrac{1}{R}\Bigl(\dfrac{1}{C_1}+\dfrac{1}{C_2}\Bigr)=\dfrac{2}{RC_1}=\dfrac{2}{\qty{20e3}{\ohm}\times\qty{100e-6}{F}}=\dfrac{2}{\qty{2}{s}}$}
なので、時定数は $\tau=\qty{2}{s}\div2=\qty{1}{s}$ です（Ω×F＝s は 6-4）。$V_1=5+5e^{-t}$、$V_2=5-5e^{-t}$〔V〕になります。

\begin{bkcenter}
\begin{bktable}{width=\linewidth, colspec={Q[l,m,wd=24mm] X[c,m] X[c,m] X[c,m]}, rowsep=1.5pt, row{1}={font=\sffamily\small}}
 & 水そう & 熱 & 電気 \\
そろう量 & 水位 $h$ & 温度 $T$ & 電圧 $V$ \\
ためる部分 & 底面積 $A$ & 熱容量 $C_{\mathrm{th}}$ & 静電容量 $C$ \\
流れにくさ & 管の流れにくさ & 熱抵抗 $R_{\mathrm{th}}$ & 抵抗 $R$ \\
一定の量 & 水の体積 $Ah$ の和 & 熱 $C_{\mathrm{th}}T$ の和 & 電気の量 $CV$ の和 \\
\end{bktable}
\bktabcaption{3つの「そろっていく」しくみは、同じ形の連立微分方程式で書ける（6-15 の相似の表へ）。}
\end{bkcenter}

\begin{taisetsu}
\def\tsl#1{\makebox[7.6\zw][l]{\sffamily\bfseries\small\color{\bkTT}#1}}
\tsl{連立微分方程式}いくつかの量の変わる速さが、おたがいの量で決まる\\[3pt]
\tsl{表計算}全部の列を、「今の行」の値で同時に進める\\[3pt]
\tsl{和と差}（重みつきの）和は一定。差は $e^{-t/\tau}$ でへる\\[3pt]
\tsl{振動モード}いっしょ $\omega_1=\sqrt{k/m}$、反対 $\omega_2=\sqrt{(k+2K)/m}$\\[3pt]
\tsl{}どんな動きも、モードのたし算\\[3pt]
\tsl{消去法}1つの文字の式にまとめて、特性方程式（6-12）
\end{taisetsu}

\begin{kotoba}
\small
\bkword{連立微分方程式}{いくつかの量の変わる速さが、おたがいの量で決まるような、微分方程式の組。}
\bkword{状態}{その時刻の位置と速度のように、その先の動きを決めるのに必要な量の組。}
\bkword{重みつきの和}{底面積や熱容量のような「重み」をかけてから、たしたもの。}
\bkword{振動モード}{全部の部分が、同じ角振動数でそろってゆれる動き方。角振動数の小さいほうから1次モード、2次モード。}
\bkword{消去法}{1つの式を「$h_2=\cdots$」の形にして、もう1つの式に入れ、文字を1つ消す方法。}
\end{kotoba}

\bkkigocards{\bkkigocard{$x_1,\ x_2$}{エックス・いち、エックス・に}{下つきの数で、どの物の量かを区別する。6-2 の $y_n$（$n$ 歩め）とはちがう}{subscript}{6-14}}
{\bkkigocard{$\omega_1,\ \omega_2$}{オメガ・いち、オメガ・に}{1次モード・2次モードの角振動数（rad/s）。6-6 の $\omega_{\mathrm n}$ の n は番号ではない}{mode frequency}{6-14}}

\begin{motto}[4]
\kviobox{電気のエネルギーは、半分になる}
図の2つのコンデンサーで、エネルギー $\tfrac12CV^2$（5-12）をくらべます。
\kvioline{はじめ：$\dfrac12\times\qty{100e-6}{F}\times(\qty{10}{V})^2=\qty{5}{mJ}$}
\kvioline{最後：$2\times\dfrac12\times\qty{100e-6}{F}\times(\qty{5}{V})^2=\qty{2.5}{mJ}$}
電気の量は保たれるのに、エネルギーはちょうど半分になりました。のこりの半分は、電流が流れるあいだに抵抗で熱になったのです。ふしぎなことに、この「半分」は抵抗の大きさによりません。抵抗が小さいと、大きな電流が短い時間流れ、熱になる量は同じになるのです。
\end{motto}

\begin{senmonsho}[大学の教科書ふうの書き方]
\begin{bkhonbun}
連立1階線形微分方程式 $\dd\symbfit{x}/\dd t=A\symbfit{x}$ において、行列 $A$ の固有値を $\lambda_1$、$\lambda_2$、対応する固有ベクトルを $\symbfit{v}_1$、$\symbfit{v}_2$ とすると、一般解は $\symbfit{x}(t)=c_1\symbfit{v}_1e^{\lambda_1t}+c_2\symbfit{v}_2e^{\lambda_2t}$ で与えられる。
\end{bkhonbun}
\bkyaku{「$\symbfit{x}$」は $h_1$、$h_2$ のように、いくつかの量をまとめて1つの文字にしたもの（第7巻のベクトル）。「行列 $A$」は⑥の係数の表。「固有ベクトル」は和や差のような、ほどける組み合わせで、「固有値」はそのときの $s$（この本では $s$、大学の本では $\lambda$ と書くことが多い）。「答えは、モードごとの指数関数のたし算」という意味。}
\end{senmonsho}

\begin{motto}[4]
\kviobox{2段の RC：なめらかにする回路}
抵抗とコンデンサーの組（6-4 の RC）を2段につなぐと、連立微分方程式になります。電池 $E=\qty{10}{V}$、どちらも $R=\qty{20}{k\ohm}$、$C=\qty{100}{\micro F}$（$RC=\qty{2}{s}$）のとき、2つの電圧は
\kvioline{$\dfrac{\dd V_1}{\dd t}=\dfrac{(E-V_1)-(V_1-V_2)}{RC}$\qquad $\dfrac{\dd V_2}{\dd t}=\dfrac{V_1-V_2}{RC}$}
1つめのコンデンサーには、電池から入る電流 $(E-V_1)/R$ と、2つめへ出ていく電流 $(V_1-V_2)/R$ の差がたまるからです。表計算で計算すると、2段めの電圧 $V_2$ は、はじめは平らで、ゆっくり立ち上がります（図）。急な変化をなめらかにするので、電気の雑音を消す「フィルター」に使われます。
\begin{bkcenter}
\begin{tikzpicture}
\begin{axis}[bkplot, width=86mm, height=40mm, xmin=0, xmax=16.8, ymin=0, ymax=11.5, xtick={4,8,12,16}, ytick={5,10}, xlabel={$t$（s）}, ylabel={電圧（V）}]
  \addplot[bkthin, dashed, domain=0:16.3] {10};
  \addplot[bkcurve, domain=0:16.3, samples=60] {10*(1-exp(-x/2))};
  \addplot[bkcurve2] coordinates {\kviorctwo};
  % Ver2 No.37：色をつけた字（━）をやめ、線の見本にした
  \node[font=\sffamily\footnotesize, anchor=south east, align=left, inner sep=1.5pt] at (axis cs:16.3,0.6) {\tikz[baseline=-0.6ex]\draw[bkBlue, line width=1.7pt] (0,0) -- (5mm,0);\ 1段の RC\\\tikz[baseline=-0.6ex]\draw[bkOrange, line width=1.7pt] (0,0) -- (5mm,0);\ 2段めの $V_2$};
\end{axis}
\end{tikzpicture}
\end{bkcenter}
\end{motto}

\begin{motto}[4]
\kviobox{えさと、食べる動物の数}
ウサギ（えさ）の数 $x$ とキツネ（食べる動物）の数 $y$ の変わり方を、$\dot x=x-0.05xy$、$\dot y=-y+0.01xy$ と表すことがあります（ロトカ1925年、ヴォルテラ1926年）。$xy$ は「出会う回数」を表します。量どうしのかけ算が入るので線形ではなく（5-18）、式では解けませんが、表計算なら解けます。ウサギがふえると少しおくれてキツネがふえ、ウサギがへるとキツネもへる、というくり返しになります。
\end{motto}

\begin{column}{カオス：未来は、はじめの状態で決まる？}
1961年、アメリカの気象学者ローレンツは、12個の量の連立微分方程式で、コンピュータに天気の動きを計算させていました。ある日、計算をとちゅうからやり直そうとして、前の計算結果の0.506127という数を、印刷されていた0.506で入れました。ほんの0.000127のちがいです。ところが、しばらくすると、天気の動きはまったくちがうものになってしまいました。

ローレンツは、1963年の論文で、これを3つの量 $x$、$y$、$z$ の式にまで簡単にしてみせました。
\kvioline{$\dot x=10(y-x)$\qquad $\dot y=x(28-z)-y$\qquad $\dot z=xy-\tfrac{8}{3}z$}
$xz$ や $xy$ のかけ算が入っているので、線形ではありません。表計算で解くと、状態の点は、チョウの羽のような形の上を、左の輪と右の輪をでたらめにうつりながら、いつまでも回り続けます（$x$ と $z$ のグラフ）。はじめの $y$ を1と1.001にした2つの計算は、しばらく同じ動きをしますが、18秒ほどで、まったくちがう動きになります（青とオレンジの線）。
\begin{bkcenter}
\begin{tikzpicture}
\begin{axis}[bkplot, axis lines=box, width=46mm, height=44mm, xmin=-21, xmax=21, ymin=0, ymax=50, xtick={-15,0,15}, ytick={10,25,40}, xlabel={$x$}, ylabel={$z$}, every axis x label/.style={at={(axis description cs:0.5,0)}, anchor=north, yshift=-5mm}, every axis y label/.style={at={(axis description cs:0,1)}, anchor=south, yshift=1mm}, grid=none]
  \addplot[bkBlue, line width=0.35pt] coordinates {\kviolorenzxz};
\end{axis}
\end{tikzpicture}\hspace{4mm}%
\begin{tikzpicture}
\begin{axis}[bkplot, axis lines=box, grid=none, width=62mm, height=44mm, xmin=0, xmax=31, ymin=-22, ymax=22, xtick={0,10,20,30}, ytick={-15,0,15}, xlabel={$t$}, ylabel={$x$}, every axis x label/.style={at={(axis description cs:0.5,0)}, anchor=north, yshift=-5mm}, every axis y label/.style={at={(axis description cs:0,1)}, anchor=south, yshift=1mm}]
  \addplot[bkBlue, line width=0.9pt] coordinates {\kviolorenza};
  \addplot[bkOrange, line width=0.9pt] coordinates {\kviolorenzb};
\end{axis}
\end{tikzpicture}
\end{bkcenter}
ほんの少しのちがいが、どんどん大きく育つ性質を、\textbf{カオス}といいます。ルールがはっきり決まっていても、はじめの状態を「完全に」は知ることができないので、遠い先までは予想できないのです。天気予報が1〜2週間より先を当てられないわけの1つが、これです。1972年の学会でのローレンツの講演の題「ブラジルでの1匹のチョウのはばたきは、テキサスで竜巻を起こすか」から、この性質は「バタフライ効果」ともよばれます。
\end{column}

%% ---------------- 例題 ----------------
\bksection{reidai}{例題でたしかめよう}

\begin{reidai}[表計算の1行を手で]
同じ形の2つの水そうの水位 $h_1$、$h_2$〔cm〕が、次の式にしたがって変わります（$t$ は分）。
\kvioline{$\dfrac{\dd h_1}{\dd t}=-0.2(h_1-h_2)$\qquad $\dfrac{\dd h_2}{\dd t}=0.2(h_1-h_2)$}
はじめ $h_1=30$、$h_2=10$ です。$\Delta t=0.5$ 分のオイラー法で、0.5分後と1分後の水位を求めましょう。\par
\kangaekata ③のように、2つとも「今の行」の差を使って進める。合計が変わらないことで確かめる。\par
\shiki\kvios\hangindent=2.6\zw\hangafter=1 0.5分後：今の差は $30-10=20$、変わる量は $0.2\times20\times0.5=2$\par
\kviob $h_1=30-2=28$、\ $h_2=10+2=12$\par
\kvioa 1分後：今の差は $28-12=16$、変わる量は $0.2\times16\times0.5=1.6$\par
\kviob $h_1=28-1.6=26.4$、\ $h_2=12+1.6=13.6$\par
\kvioa 確かめ：合計は $28+12=40$、\ $26.4+13.6=40$。はじめの $30+10=40$ と同じ。\par
\kotae\kvios\hangindent=2.6\zw\hangafter=1 0.5分後 $h_1=\qty{28}{cm}$、$h_2=\qty{12}{cm}$\par
\kvioa 1分後 $h_1=\qty{26.4}{cm}$、$h_2=\qty{13.6}{cm}$
\end{reidai}

\begin{jibun}
次の式で、はじめ $h_1=25$、$h_2=5$（cm）です。$\Delta t=1$ 分のオイラー法で、1分後と2分後の水位を求めましょう（$t$ は分）。
\kvioline{$\dfrac{\dd h_1}{\dd t}=-0.1(h_1-h_2)$\qquad $\dfrac{\dd h_2}{\dd t}=0.1(h_1-h_2)$}
\bkjans{1分後：差20、変わる量 $0.1\times20\times1=2$ で、$h_1=\qty{23}{cm}$、$h_2=\qty{7}{cm}$。2分後：差16、変わる量1.6で、$h_1=\qty{21.4}{cm}$、$h_2=\qty{8.6}{cm}$。}
\end{jibun}

\begin{reidai}[2回微分の式を、1回微分の式2つにする]
ダンパーつきのばね $0.5\ddot x+0.4\dot x+8x=0$ を、位置 $x$ と速度 $v$ の1回微分の式2つに書きかえましょう。また、$x=0.1$、$v=0$ から、$\Delta t=0.1$ の「速度を先に進めるオイラー法」で2行進めましょう。\par
\kangaekata $\dot x=v$、$\ddot x=\dot v$（5-5）。もとの式を「$\ddot x=\cdots$」の形にする。表は6-2の順（加速度→速度→位置）。\par
\shiki\kvios\hangindent=2.6\zw\hangafter=1 両辺を0.5でわる：$\ddot x=-16x-0.8\dot x$\par
\kvioa 1回微分の式2つ：$\dot x=v$、\ $\dot v=-16x-0.8v$\par
\kvioa 1行め：$a=-16\times0.1-0.8\times0=-1.6$\par
\kviob $v=0+(-1.6)\times0.1=-0.16$\par
\kviob $x=0.1+(-0.16)\times0.1=0.084$\par
\kvioa 2行め：$a=-16\times0.084-0.8\times(-0.16)$\par
\kviob $=-1.344+0.128=-1.216$\par
\kviob $v=-0.16+(-1.216)\times0.1=-0.2816$\par
\kviob $x=0.084+(-0.2816)\times0.1=0.05584$\par
\kotae\kvios\hangindent=2.6\zw\hangafter=1 $\dot x=v$、$\dot v=-16x-0.8v$\par
\kvioa $(x,\ v)$ は $(0.084,\ -0.16)$、$(0.05584,\ -0.2816)$
\end{reidai}

\begin{jibun}
$\ddot x+0.2\dot x+9x=0$ を、1回微分の式2つに書きかえましょう。また、$x=0.2$、$v=0$ から、$\Delta t=0.1$ の「速度を先に進めるオイラー法」で2行進めましょう。
\bkjans{$\dot x=v$、$\dot v=-9x-0.2v$。1行め：$a=-1.8$、$v=-0.18$、$x=0.182$。2行め：$a=-9\times0.182-0.2\times(-0.18)=-1.602$、$v=-0.3402$、$x=0.14798$。}
\end{jibun}

\begin{reidai}[和と差で解く]
同じ形の2つの水そうの水位が、次の式にしたがいます（cm、分）。
\kvioline{$\dfrac{\dd h_1}{\dd t}=-0.2(h_1-h_2)$\qquad $\dfrac{\dd h_2}{\dd t}=0.2(h_1-h_2)$}
はじめ $h_1=30$、$h_2=10$ です。$h_1$、$h_2$ を $t$ の式で表しましょう。また、2つの水位の差が\qty{10}{cm}になるのは何分後ですか（$\ln2=0.693$）。\par
\kangaekata ⑤のように、和 $S$ と差 $D$ の式を作る。$S$ は一定、$D$ は6-4の形。最後に $h_1=(S+D)\div2$、$h_2=(S-D)\div2$ にもどす。\par
\shiki\kvios\hangindent=2.6\zw\hangafter=1 たす：$\dfrac{\dd S}{\dd t}=0$ なので $S=30+10=40$\par
\kvioa ひく：$\dfrac{\dd D}{\dd t}=-0.2D-0.2D=-0.4D$\par
\kviob はじめ $D=30-10=20$ なので $D=20e^{-0.4t}$\par
\kvioa $h_1=\dfrac{40+20e^{-0.4t}}{2}=20+10e^{-0.4t}$\par
\kvioa $h_2=\dfrac{40-20e^{-0.4t}}{2}=20-10e^{-0.4t}$\par
\kvioa 差が10：$20e^{-0.4t}=10$ より $e^{-0.4t}=\dfrac12$\par
\kviob 両辺の $\ln$ をとる：$-0.4t=-\ln2$\par
\kviob $t=\dfrac{0.693}{0.4}\approx1.73$\par
\kotae\kvios\hangindent=2.6\zw\hangafter=1 $h_1=20+10e^{-0.4t}$、$h_2=20-10e^{-0.4t}$。約1.7分後
\end{reidai}

\begin{jibun}
同じ形の2つの水そうで、次の式にしたがい（cm、分）、はじめ $h_1=40$、$h_2=0$ です。$h_1$、$h_2$ を $t$ の式で表し、差が\qty{20}{cm}になる時刻を求めましょう（$\ln2=0.693$）。
\kvioline{$\dfrac{\dd h_1}{\dd t}=-0.05(h_1-h_2)$\qquad $\dfrac{\dd h_2}{\dd t}=0.05(h_1-h_2)$}
\bkjans{$S=40$。$\dfrac{\dd D}{\dd t}=-0.1D$、$D=40e^{-0.1t}$。$h_1=20+20e^{-0.1t}$、$h_2=20-20e^{-0.1t}$。$e^{-0.1t}=\dfrac12$ より $t=6.93$ で、約6.9分後。}
\end{jibun}

\begin{reidai}[2つの物の温度]
80\,℃のお湯200\,g（熱容量 $C_1=\qty{840}{J/K}$）を入れた金属のコップを、20\,℃の水600\,g（$C_2=\qty{2520}{J/K}$）にひたします。熱抵抗は $R_{\mathrm{th}}=\qty{0.5}{K/W}$ で、外へ熱は逃げません。最後の温度、差がへるときの時定数、2つの温度の式を求めましょう。\par
\kangaekata ⑧の式。熱容量をかけた和（熱）は一定。差 $D=T_1-T_2$ は、6-4 の形でへる。\par
\shiki\kvios\hangindent=2.6\zw\hangafter=1 最後の温度：$T_\infty=\dfrac{840\times80+2520\times20}{840+2520}=\dfrac{\num{117600}}{3360}=35$\par
\kvioa $\dfrac{1}{C_1}+\dfrac{1}{C_2}=\dfrac{3}{2520}+\dfrac{1}{2520}=\dfrac{4}{2520}=\dfrac{1}{630}$\par
\kvioa $\dfrac{1}{\tau}=\dfrac{1}{0.5}\times\dfrac{1}{630}=\dfrac{1}{315}$ なので $\tau=\qty{315}{s}$\par
\kvioa 差：はじめ $80-20=60$ なので $D=60e^{-t/315}$\par
\kvioa 熱は一定：$840T_1+2520T_2=3360\times35$\par
\kvioa $T_2=T_1-D$ を入れる：$3360T_1-2520D=3360\times35$\par
\kviob $T_1=35+\dfrac{2520}{3360}D=35+\dfrac{3}{4}D=35+45e^{-t/315}$\par
\kviob $T_2=T_1-D=35-15e^{-t/315}$\par
\kotae\kvios\hangindent=2.6\zw\hangafter=1 35\,℃。$\tau=\qty{315}{s}$（約5分）。\par
\kvioa $T_1=35+45e^{-t/315}$、$T_2=35-15e^{-t/315}$
\end{reidai}
\begin{jibun}
90\,℃のお湯100\,g（$C_1=\qty{420}{J/K}$）のコップを、10\,℃の水300\,g（$C_2=\qty{1260}{J/K}$）にひたします。$R_{\mathrm{th}}=\qty{1}{K/W}$ で、外へ熱は逃げません。最後の温度、時定数、2つの温度の式を求めましょう。
\bkjans{$T_\infty=\dfrac{420\times90+1260\times10}{1680}=30$\,℃。$\dfrac{1}{C_1}+\dfrac{1}{C_2}=\dfrac{4}{1260}=\dfrac{1}{315}$ で $\tau=\qty{315}{s}$。$D=80e^{-t/315}$、$T_1=30+60e^{-t/315}$、$T_2=30-20e^{-t/315}$。}
\end{jibun}



\begin{reidai}[2つのコンデンサー]
\qty{10}{V}に充電した $C_1=\qty{100}{\micro F}$ のコンデンサーと、空の $C_2=\qty{400}{\micro F}$ のコンデンサーを、$R=\qty{10}{k\ohm}$ の抵抗でつなぎます。最後の電圧、時定数、$V_1$、$V_2$ の式を求めましょう。\par
\kangaekata ⑧の式。電気の量 $C_1V_1+C_2V_2$ は一定。差の時定数は、例題4の $R_{\mathrm{th}}$ を $R$ に、熱容量を静電容量にかえて求める。\par
\shiki\kvios\hangindent=2.6\zw\hangafter=1 電気の量（\unit{\micro C}）：$100\times10+400\times0=1000$。最後は $500V=1000$ で $V=2$\par
\kvioa $\dfrac{1}{C_1}+\dfrac{1}{C_2}=\dfrac{1}{\qty{100}{\micro F}}+\dfrac{1}{\qty{400}{\micro F}}=\dfrac{5}{\qty{400}{\micro F}}=\dfrac{1}{\qty{80}{\micro F}}$\par
\kvioa $\tau=R\times\qty{80}{\micro F}=\qty{10e3}{\ohm}\times\qty{80e-6}{F}=\qty{0.8}{s}$。差は $D=10e^{-t/0.8}$\par
\kvioa 電気の量の式に $V_1=V_2+D$ を入れる：$100(V_2+D)+400V_2=1000$\par
\kviob $V_2=2-\dfrac{D}{5}=2-2e^{-t/0.8}$、\ $V_1=V_2+D=2+8e^{-t/0.8}$\par
\kotae\kvios\hangindent=2.6\zw\hangafter=1 \qty{2}{V}。$\tau=\qty{0.8}{s}$。\par
\kvioa $V_1=2+8e^{-t/0.8}$、$V_2=2-2e^{-t/0.8}$（V）
\end{reidai}

\begin{jibun}
\qty{9}{V}に充電した $C_1=\qty{200}{\micro F}$ と、空の $C_2=\qty{100}{\micro F}$ を、$R=\qty{15}{k\ohm}$ でつなぎます。最後の電圧、時定数、$V_1$、$V_2$ の式を求めましょう。
\bkjans{電気の量 $200\times9=1800$（\unit{\micro C}）、最後は $300V=1800$ で\qty{6}{V}。$\dfrac{1}{C_1}+\dfrac{1}{C_2}=\dfrac{3}{\qty{200}{\micro F}}$ で $\tau=\qty{15e3}{\ohm}\times\dfrac{\qty{200}{\micro F}}{3}=\qty{1}{s}$。$D=9e^{-t}$、$V_2=6-6e^{-t}$、$V_1=6+3e^{-t}$。}
\end{jibun}



\begin{reidai}[振動モード]
$m=\qty{2}{kg}$ の台車2つを、$k=\qty{200}{N/m}$ の外側のばね2本と、$K=\qty{100}{N/m}$ のつなぎのばねでつなぎます。\par
（1）1次モードと2次モードの角振動数と周期を求めましょう。\par
（2）$x_1=0.05$、$x_2=0.01$（m）だけずらして、静かにはなしました。$x_1$、$x_2$ を求めましょう。\par
\kangaekata ⑥の式。和 $u$、差 $w$ は、はじめの値を大きさとする $\cos$。\par
\shiki\kvios\hangindent=4.3\zw\hangafter=1 （1）$\omega_1=\sqrt{\dfrac{200}{2}}=10$、周期 $\dfrac{2\pi}{10}\approx0.63$\par
\kviob $\omega_2=\sqrt{\dfrac{200+200}{2}}=\sqrt{200}\approx14.1$、周期 $\dfrac{2\pi}{14.1}\approx0.44$\par
\kvioq （2）$u=0.05+0.01=0.06$、$w=0.05-0.01=0.04$ から、$u=0.06\cos10t$、$w=0.04\cos14.1t$\par
\kviob $x_1=\dfrac{u+w}{2}=0.03\cos10t+0.02\cos14.1t$\par
\kviob $x_2=\dfrac{u-w}{2}=0.03\cos10t-0.02\cos14.1t$\par
\kotae\kvios\hangindent=4.3\zw\hangafter=1 （1）\qty{10}{rad/s}（約\qty{0.63}{s}）、約\qty{14.1}{rad/s}（約\qty{0.44}{s}）\quad （2）式の段の $x_1$、$x_2$
\end{reidai}

\begin{jibun}
$m=\qty{1}{kg}$ の台車2つを、$k=\qty{16}{N/m}$ のばね2本と $K=\qty{10}{N/m}$ のつなぎのばねでつなぎます。\newline（1）$\omega_1$、$\omega_2$ を求めましょう。\newline（2）$x_1=0$、$x_2=0.06$（m）だけずらして静かにはなしたときの $x_1$、$x_2$ を求めましょう。
\bkjans{（1）\nobreak $\omega_1=\sqrt{16}=\qty{4}{rad/s}$、$\omega_2=\sqrt{36}=\qty{6}{rad/s}$。（2）\nobreak $u=0.06\cos4t$、$w=-0.06\cos6t$。$x_1=0.03\cos4t-0.03\cos6t$、$x_2=0.03\cos4t+0.03\cos6t$。}
\end{jibun}



\begin{reidai}[消去法：直列の水そう]
上の水そうから下の水そうへ水が落ち、下の水そうから外へ出ていきます（cm、分）。
\kvioline{$\dfrac{\dd h_1}{\dd t}=-\dfrac{h_1}{4}$\qquad $\dfrac{\dd h_2}{\dd t}=\dfrac{h_1}{4}-\dfrac{h_2}{4}$}
はじめ $h_1=40$、$h_2=0$ です。下の水位 $h_2$ を求め、いちばん高くなる時刻と高さを求めましょう（$e^{-1}=0.368$）。\par
\kangaekata ⑦の消去法。特性方程式が重解なので、答えは $(A+Bt)e^{st}$ の形。\par
\shiki\kvios\hangindent=2.6\zw\hangafter=1 下の式に4をかけて：$h_1=4\dot h_2+h_2$\par
\kvioa 上の式 $\dot h_1=-h_1/4$ に入れる：$4\ddot h_2+\dot h_2=-\dfrac{4\dot h_2+h_2}{4}$\par
\kviob 両辺に4をかけて整理：$16\ddot h_2+8\dot h_2+h_2=0$\par
\kvioa 特性方程式 $16s^2+8s+1=(4s+1)^2=0$ より $s=-\tfrac14$（重解）、$h_2=(A+Bt)e^{-t/4}$\par
\kvioa $h_2(0)=A=0$。$\dot h_2(0)=40\div4-0=10$ で、$Bte^{-t/4}$ の微分の $t=0$ の値は $B$ なので $B=10$\par
\kvioa いちばん高い時刻：$\dot h_2=10\Bigl(1-\dfrac{t}{4}\Bigr)e^{-t/4}=0$ より $t=4$。$h_2=40e^{-1}\approx14.7$\par
\kotae\kvios\hangindent=2.6\zw\hangafter=1 $h_2=10te^{-t/4}$。4分後に約\qty{14.7}{cm}
\end{reidai}

\begin{jibun}
次の式で、はじめ $h_1=30$、$h_2=0$ です（cm、分）。$h_2$ を求め、いちばん高くなる時刻と高さを求めましょう（$e^{-1}=0.368$）。
\kvioline{$\dfrac{\dd h_1}{\dd t}=-\dfrac{h_1}{2}$\qquad $\dfrac{\dd h_2}{\dd t}=\dfrac{h_1}{2}-\dfrac{h_2}{2}$}
\bkjans{$h_1=2\dot h_2+h_2$ を上の式に入れて $4\ddot h_2+4\dot h_2+h_2=0$、$(2s+1)^2=0$、$s=-\dfrac12$。$h_2=(A+Bt)e^{-t/2}$、$A=0$、$B=\dot h_2(0)=15$。$h_2=15te^{-t/2}$。2分後に $30\times0.368\approx\qty{11.0}{cm}$。}
\end{jibun}

\begin{reidai}[親と子]
放射性物質の親の数 $N_1$ と子の数 $N_2$ が、次の式にしたがいます（$t$ は時間〔h〕）。
\kvioline{$\dfrac{\dd N_1}{\dd t}=-\dfrac{N_1}{2}$\qquad $\dfrac{\dd N_2}{\dd t}=\dfrac{N_1}{2}-N_2$}
はじめ $N_1=1000$、$N_2=0$ です。$N_2$ を求め、子がいちばん多くなる時刻と数を求めましょう（$\ln2=0.693$）。\par
\kangaekata 消去法。子の式を $N_1=\cdots$ にして親の式に入れる。\par
\shiki\kvios\hangindent=2.6\zw\hangafter=1 子の式に2をかけて：$N_1=2(\dot N_2+N_2)$\par
\kvioa 親の式に入れる：$2(\ddot N_2+\dot N_2)=-(\dot N_2+N_2)$\par
\kvioa 整理：$2\ddot N_2+3\dot N_2+N_2=0$。特性方程式は $(2s+1)(s+1)=0$\par
\kviob $s=-\dfrac12,\ -1$ なので $N_2=Ae^{-t/2}+Be^{-t}$\par
\kvioa $N_2(0)=A+B=0$ より $B=-A$。$\dot N_2(0)=1000\div2-0=500$ より $-\dfrac{A}{2}-B=500$\par
\kviob $B=-A$ を入れて $\dfrac{A}{2}=500$、$A=1000$、$B=-1000$\par
\kvioa $\dot N_2=1000\Bigl(-\dfrac12e^{-t/2}+e^{-t}\Bigr)=0$ より $e^{-t}=\dfrac12e^{-t/2}$\par
\kviob 両辺を $e^{-t/2}$ でわって $e^{-t/2}=\dfrac12$、$t=2\ln2\approx1.39$\par
\kviob $e^{-t}=(e^{-t/2})^2=\dfrac14$ なので、$N_2=1000\Bigl(\dfrac12-\dfrac14\Bigr)=250$\par
\kotae\kvios\hangindent=2.6\zw\hangafter=1 $N_2=1000(e^{-t/2}-e^{-t})$。約1.4時間後に250
\end{reidai}

\begin{jibun}
飲み薬が胃腸から血液に入り、血液から体の外へ出ていきます。胃腸の薬の量 $G$、血液の薬の量 $B$〔mg〕が次の式にしたがい（$t$ は時間〔h〕）、はじめ $G=200$、$B=0$ です。$B$ を求め、いちばん多くなる時刻と量を求めましょう（$\ln2=0.693$）。
\kvioline{$\dfrac{\dd G}{\dd t}=-G$\qquad $\dfrac{\dd B}{\dd t}=G-\dfrac{B}{2}$}
\bkjans{$G=\dot B+\dfrac B2$ を入れて $\ddot B+\dfrac32\dot B+\dfrac12B=0$、$2s^2+3s+1=0$、$s=-\dfrac12,\ -1$。$B=Ae^{-t/2}+Ce^{-t}$、$A+C=0$、$\dot B(0)=200$ より $\dfrac A2=200$、$A=400$。$B=400(e^{-t/2}-e^{-t})$。$e^{-t/2}=\dfrac12$ で、$t=2\ln2\approx1.4$ 時間後に\qty{100}{mg}。}
\end{jibun}

%% ---------------- まちがえやすいところ ----------------
\bksection{tsumazuki}{まちがえやすいところをチェック}

\begin{tsumazuki}{① 片方を先に新しくしてしまう}
\bkng{先に出した $h_1=18$ を使って $h_2=0+0.1\times(18-0)=1.8$}
\bkok{今の行の差20を使って $h_2=0+0.1\times20=2$}
×では合計が $18+1.8=19.8$ になり、水が消えてしまいます。左から出た水と右に入る水は、同じ時刻の差で決まります。
\end{tsumazuki}

\begin{tsumazuki}{② 和と差のまま答える}
\bkng{$h_1=20e^{-0.2t}$}
\bkok{$h_1=10+10e^{-0.2t}$}
$20e^{-0.2t}$ は差 $D$ です。最後に $h_1=(S+D)\div2$ にもどすのを忘れないようにしましょう。
\end{tsumazuki}

\begin{tsumazuki}{③ 代入するとき微分を忘れる}
\bkng{$h_2=h_1+10\dot h_1$ を微分して、$\dot h_2=\dot h_1+10\dot h_1$}
\bkok{$\dot h_2=\dot h_1+10\ddot h_1$}
$10\dot h_1$ を微分すると $10\ddot h_1$（2回微分）になります。
\end{tsumazuki}

\noindent\begin{minipage}[t]{0.485\linewidth}\vspace{0pt}
\begin{tsumazuki}{④ モードの $\omega$ をたす}
\raggedright\bkng{$x_1$ は $\omega_1+\omega_2$ でゆれる}
\bkok{$\omega_1$ のゆれと $\omega_2$ のゆれのたし算}
たすのは $\cos$ の式どうしで、角振動数どうしではありません。
\end{tsumazuki}
\end{minipage}\hfill
\begin{minipage}[t]{0.485\linewidth}\vspace{0pt}
\begin{tsumazuki}{⑤ 重みを忘れる}
\raggedright\bkng{最後は $20\div2=\qty{10}{cm}$}
\bkok{体積 $2000\div400=\qty{5}{cm}$}
100\,cm² と300\,cm² の水そうで一定なのは、水の体積です。
\end{tsumazuki}
\end{minipage}

%% ---------------- 物理・工学ではこう使う ----------------
\begin{butsuri}
\kviobox{飲み薬の血液中の量}
\noindent\begin{minipage}[c]{0.54\linewidth}
\raggedright
薬は胃腸から血液に入り、肝臓やじん臓で取りのぞかれます。血液の薬の量は、じぶんでとこう8と同じ形の式で、ふえてからへります。何時間おきに飲むかは、このような式と、たくさんの人の測定で決められています（薬物動態学）。
\end{minipage}\hfill
\begin{minipage}[c]{0.42\linewidth}
\centering
\begin{tikzpicture}
\begin{axis}[bkplot, width=46mm, height=32mm, xmin=0, xmax=8.6, ymin=0, ymax=125, xtick={2,4,6,8}, ytick={50,100}, xlabel={$t$（h）}, ylabel={mg}]
  \addplot[bkGreen!70!black, line width=0.9pt, dashed, domain=0:8.3] {50};
  \addplot[bkcurve2, domain=0:8.3, samples=60] {400*(exp(-x/2)-exp(-x))};
  \node[font=\sffamily\footnotesize, anchor=south west, text=bkGreen!60!black] at (axis cs:4.6,51) {効く量（例）};
\end{axis}
\end{tikzpicture}
\end{minipage}
\end{butsuri}

\begin{butsuri}
\kviobox{建物の各階のゆれ：振動モード}
\noindent\begin{minipage}[c]{0.58\linewidth}
\raggedright
建物は、各階の床（おもり）が、柱とかべ（ばね）でつながったものと考えられます。2階建てなら⑥の台車と同じ形の式になり、1次モード（全部の階が同じ向き）と2次モード（上と下が反対）があります。地震のゆれの周期がモードの周期に近いと、大きくゆれます（6-13 の共振）。1次モードの周期は、おおよそ「階数×0.1秒」といわれます（目安）。
\end{minipage}\hfill
\begin{minipage}[c]{0.38\linewidth}
\centering
\begin{tikzpicture}[x=0.9mm, y=0.75mm]
  \fill[bkInk!50] (-2,-1.2) rectangle (52,0);
  \draw[bkGray, line width=0.4pt, dash pattern=on 1.5pt off 1.5pt] (0,0) -- (0,30) (16,0) -- (16,30) (32,0) -- (32,30) (48,0) -- (48,30);
  \draw[bkInk!70, line width=1pt] (0,0) .. controls (0,6) and (3,9) .. (3,14) .. controls (3,20) and (5,23) .. (5,28);
  \draw[bkInk!70, line width=1pt] (16,0) .. controls (16,6) and (19,9) .. (19,14) .. controls (19,20) and (21,23) .. (21,28);
  \filldraw[fill=bkBlue!30, draw=bkBlue!80!black, line width=0.7pt] (2,13) rectangle (20,15);
  \filldraw[fill=bkBlue!30, draw=bkBlue!80!black, line width=0.7pt] (4,27) rectangle (22,29);
  \draw[bkInk!70, line width=1pt] (32,0) .. controls (32,6) and (35,9) .. (35,14) .. controls (35,20) and (31,23) .. (30,28);
  \draw[bkInk!70, line width=1pt] (48,0) .. controls (48,6) and (51,9) .. (51,14) .. controls (51,20) and (47,23) .. (46,28);
  \filldraw[fill=bkOrange!30, draw=bkOrange!80!black, line width=0.7pt] (34,13) rectangle (52,15);
  \filldraw[fill=bkOrange!30, draw=bkOrange!80!black, line width=0.7pt] (29,27) rectangle (47,29);
  \node[font=\sffamily\footnotesize, anchor=north] at (11,-2.5) {1次モード};
  \node[font=\sffamily\footnotesize, anchor=north] at (41,-2.5) {2次モード};
\end{tikzpicture}
\end{minipage}
\end{butsuri}

\begin{butsuri}
\kviobox{病院の「テクネチウム・ジェネレーター」：親と子}
\noindent\begin{minipage}[c]{0.56\linewidth}
\raggedright
体の中を写す検査に使うテクネチウム99m（半減期約6時間）は、作りおきができません。そこで病院では、親のモリブデン99（半減期約66時間）の容器をおき、生まれた子を必要なときに取り出します。子は例題8と同じ形の式でふえ、取り出してから約23時間後にいちばん多くなります（グラフは1時間にこわれる数で、はじめの親を1とした。親の約9割が子になる）。
\end{minipage}\hfill
\begin{minipage}[c]{0.4\linewidth}
\centering
\begin{tikzpicture}
\begin{axis}[bkplot, width=44mm, height=40mm, xmin=0, xmax=76, ymin=0, ymax=1.15, xtick={23,48,72}, ytick={0.5,1}, xlabel={$t$（h）}, yticklabel style={/pgf/number format/fixed}]
  \addplot[bkthin, domain=0:72, samples=40] {exp(-0.010512*x)};
  \addplot[bkcurve2, domain=0:72, samples=80] {(0.11533/(0.11533-0.010512))*(exp(-0.010512*x)-exp(-0.11533*x))};
  \node[font=\sffamily\footnotesize, anchor=south, text=bkGray] at (axis cs:60,0.6) {親};
  \node[font=\sffamily\footnotesize, anchor=north, text=bkOrange!80!black] at (axis cs:40,0.62) {子};
\end{axis}
\end{tikzpicture}
\end{minipage}
\end{butsuri}


%% ---------------- 章末問題 ----------------
\bkendproblems
\bkproblemgroup{基本}{5}

\bkproblem $\dfrac{\dd h_1}{\dd t}=-0.3(h_1-h_2)$、$\dfrac{\dd h_2}{\dd t}=0.3(h_1-h_2)$（cm、分）、はじめ $h_1=50$、$h_2=10$ です。$\Delta t=0.5$ 分のオイラー法で、0.5分後と1分後の水位を求めましょう\mbox{（\bkreidairef{1}）}。
\bkpans{0.5分後：差40、変わる量 $0.3\times40\times0.5=6$ で $h_1=\qty{44}{cm}$、$h_2=\qty{16}{cm}$。1分後：差28、変わる量4.2で $h_1=\qty{39.8}{cm}$、$h_2=\qty{20.2}{cm}$。}

\bkproblem $\ddot x+0.5\dot x+25x=0$ を1回微分の式2つに書きかえ、$x=0.1$、$v=0$ から、$\Delta t=0.1$ の「速度を先に進めるオイラー法」で1行進めましょう\mbox{（\bkreidairef{2}）}。
\bkpans{$\dot x=v$、$\dot v=-25x-0.5v$。$a=-2.5$、$v=-0.25$、$x=0.1+(-0.25)\times0.1=0.075$。}

\bkproblem 同じ形の2つの水そうで、$\dfrac{\dd h_1}{\dd t}=-0.1(h_1-h_2)$、$\dfrac{\dd h_2}{\dd t}=0.1(h_1-h_2)$（cm、分）、はじめ $h_1=12$、$h_2=4$ です。$h_1$、$h_2$ を求め、差が半分の\qty{4}{cm}になる時刻を求めましょう\mbox{（\bkreidairef{3}）}（$\ln2=0.693$）。
\bkpans{$S=16$、$D=8e^{-0.2t}$。$h_1=8+4e^{-0.2t}$、$h_2=8-4e^{-0.2t}$。$e^{-0.2t}=\dfrac12$ より $t=5\ln2\approx3.5$ 分後。}

\bkproblem 60\,℃のお湯200\,g（$C_1=\qty{840}{J/K}$）のコップを、15\,℃の水400\,g（$C_2=\qty{1680}{J/K}$）にひたします（$R_{\mathrm{th}}=\qty{0.5}{K/W}$、外へ熱は逃げない）。最後の温度、時定数、2つの温度の式を求めましょう\mbox{（\bkreidairef{4}）}。
\bkpans{$T_\infty=\dfrac{840\times60+1680\times15}{2520}=30$\,℃。$\dfrac{1}{C_1}+\dfrac{1}{C_2}=\dfrac{3}{1680}$ で $\tau=0.5\times560=\qty{280}{s}$。$T_1=30+30e^{-t/280}$、$T_2=30-15e^{-t/280}$。}

\bkproblem $m=\qty{0.5}{kg}$ の台車2つを、$k=\qty{50}{N/m}$ のばね2本と $K=\qty{25}{N/m}$ のつなぎのばねでつなぎます。\newline（1）$\omega_1$、$\omega_2$ を求めましょう。\newline（2）$x_1=0.04$、$x_2=0$（m）だけずらして静かにはなしたときの $x_1$、$x_2$ を求めましょう\mbox{（\bkreidairef{6}）}。
\bkpans{（1）\nobreak $\omega_1=\sqrt{100}=\qty{10}{rad/s}$、$\omega_2=\sqrt{100\div0.5}=\sqrt{200}\approx\qty{14.1}{rad/s}$。（2）\nobreak $u=0.04\cos10t$、$w=0.04\cos14.1t$。$x_1=0.02\cos10t+0.02\cos14.1t$、$x_2=0.02\cos10t-0.02\cos14.1t$。}

\bkproblem $\dfrac{\dd h_1}{\dd t}=-\dfrac{h_1}{5}$、$\dfrac{\dd h_2}{\dd t}=\dfrac{h_1}{5}-\dfrac{h_2}{5}$（cm、分）、はじめ $h_1=50$、$h_2=0$ です。$h_2$ を求め、いちばん高くなる時刻と高さを求めましょう\mbox{（\bkreidairef{7}）}（$e^{-1}=0.368$）。
\bkpans{$h_1=5\dot h_2+h_2$ を入れて $25\ddot h_2+10\dot h_2+h_2=0$、$(5s+1)^2=0$。$h_2=(A+Bt)e^{-t/5}$、$A=0$、$B=10$。$h_2=10te^{-t/5}$。5分後に $50\times0.368\approx\qty{18.4}{cm}$。}

\bkproblemgroup{チャレンジ}{5}

\bkproblem \qty{6}{V}に充電した $C_1=\qty{100}{\micro F}$ と、空の $C_2=\qty{50}{\micro F}$ を、$R=\qty{30}{k\ohm}$ でつなぎます。最後の電圧、時定数、$V_1$、$V_2$ の式を求めましょう\mbox{（\bkreidairef{5}）}。
\bkpans{電気の量 \qty{600}{\micro C}、最後は $150V=600$ で\qty{4}{V}。$\dfrac{1}{C_1}+\dfrac{1}{C_2}=\dfrac{3}{\qty{100}{\micro F}}$ で $\tau=\qty{30e3}{\ohm}\times\dfrac{\qty{100}{\micro F}}{3}=\qty{1}{s}$。$D=6e^{-t}$、$V_2=4-4e^{-t}$、$V_1=4+2e^{-t}$。}

\bkproblem 親の数 $N_1$ と子の数 $N_2$ が $\dfrac{\dd N_1}{\dd t}=-\dfrac{N_1}{3}$、$\dfrac{\dd N_2}{\dd t}=\dfrac{N_1}{3}-N_2$（$t$ は時間〔h〕）にしたがい、はじめ $N_1=900$、$N_2=0$ です。$N_2$ を求め、3時間後の $N_2$ を求めましょう\mbox{（\bkreidairef{8}）}（$e^{-1}=0.368$、$e^{-3}=0.050$）。
\bkpans{$N_1=3(\dot N_2+N_2)$ を入れて $3\ddot N_2+4\dot N_2+N_2=0$、$(3s+1)(s+1)=0$。$N_2=Ae^{-t/3}+Be^{-t}$、$A+B=0$、$-\dfrac A3-B=300$ より $A=450$。$N_2=450(e^{-t/3}-e^{-t})$。3時間後 $450\times(0.368-0.050)\approx143$。}

\bkproblem 底面積\qty{200}{cm^2}の水そう（左）と\qty{100}{cm^2}の水そう（右）を、1分に $20(h_2-h_1)$\,cm³ 流れる管でつなぎます。はじめ左は空、右は\qty{30}{cm}です。最後の水位と、$h_1$、$h_2$ の式を求めましょう\mbox{（\bkreidairef{4}）}。
\bkpans{$\dfrac{\dd h_1}{\dd t}=0.1(h_2-h_1)$、$\dfrac{\dd h_2}{\dd t}=-0.2(h_2-h_1)$。体積 $200h_1+100h_2=3000$ が一定で、最後は\qty{10}{cm}。差 $D=h_2-h_1$ は $\dfrac{\dd D}{\dd t}=-0.3D$、$D=30e^{-0.3t}$。$h_1=10-10e^{-0.3t}$、$h_2=10+20e^{-0.3t}$。}

\begin{furikaeri}
\begin{checklist}
\item からみあう2つの量を、表計算で同時に進めて計算できる。\bkfurref{bk:6-14:hyou}
\item 和と差（重みつきの和と差）に分けて、2つの式を解ける。\bkfurref{bk:6-14:wasa}
\item 振動モードや消去法で、連立微分方程式の答えを求められる。\bkfurref{bk:6-14:mode}
\end{checklist}
\end{furikaeri}

\bkfukidashi{sora}{組みかえると、ばらばらに解けた！　次は第6巻のまとめだよ。}
